\documentclass[11pt,a4paper,oneside]{article}

\usepackage[T1]{fontenc}
\usepackage[utf8]{inputenc}
\usepackage{libertinus}
\usepackage{microtype}
\usepackage[a4paper,margin=27mm,headheight=15pt]{geometry}
\usepackage{amsmath,amssymb,amsthm,mathtools}
\usepackage{bm,mathrsfs}
\usepackage{booktabs,longtable,array,tabularx,multirow}
\usepackage{enumitem}
\usepackage{xcolor}
\usepackage{tikz}
\usetikzlibrary{arrows.meta,positioning,calc,decorations.pathreplacing}
\usepackage[most]{tcolorbox}
\usepackage{fancyhdr}
\usepackage{xurl}
\usepackage{hyperref}
\usepackage[nameinlink,noabbrev,capitalise]{cleveref}
\usepackage{listings}
\usepackage{etoolbox}

\definecolor{deepolive}{RGB}{67,79,45}
\definecolor{softolive}{RGB}{236,239,226}
\definecolor{darkteal}{RGB}{31,83,84}
\definecolor{warmgray}{RGB}{92,87,78}
\definecolor{softcoral}{RGB}{246,225,219}
\definecolor{linkblue}{RGB}{31,76,122}
\definecolor{codeback}{RGB}{247,247,242}

\hypersetup{
  colorlinks=true,
  linkcolor=deepolive,
  citecolor=darkteal,
  urlcolor=linkblue,
  pdftitle={The Full Asymptotic Transseries of the Fubini Numbers},
  pdfauthor={A self-contained research article with proofs},
  pdfsubject={Ordered Bell numbers, Fubini numbers, pole-lattice transseries, asymptotics, Bell polynomials}
}

\allowdisplaybreaks[3]
\raggedbottom
\emergencystretch=2em
\setcounter{tocdepth}{2}
\setcounter{secnumdepth}{3}
\setlist{itemsep=2pt,topsep=5pt,parsep=1pt}

\pagestyle{fancy}
\fancyhf{}
\fancyhead[L]{\small\nouppercase{\leftmark}}
\fancyhead[R]{\small\thepage}
\renewcommand{\headrulewidth}{0.4pt}

\newtheorem{theorem}{Theorem}[section]
\newtheorem{lemma}[theorem]{Lemma}
\newtheorem{proposition}[theorem]{Proposition}
\newtheorem{corollary}[theorem]{Corollary}
\newtheorem{identity}[theorem]{Identity}
\theoremstyle{definition}
\newtheorem{definition}[theorem]{Definition}
\newtheorem{example}[theorem]{Example}
\newtheorem{algorithm}[theorem]{Algorithm}
\theoremstyle{remark}
\newtheorem{remark}[theorem]{Remark}
\newtheorem{historical}[theorem]{Historical note}

\newtcolorbox{masterbox}[1][]{
  enhanced,breakable,colback=softolive,colframe=deepolive,
  boxrule=0.7pt,arc=1mm,left=2mm,right=2mm,top=1.5mm,bottom=1.5mm,
  title={Master result},fonttitle=\bfseries,#1
}
\newtcolorbox{interpretbox}[1][]{
  enhanced,breakable,colback=white,colframe=darkteal,
  boxrule=0.6pt,arc=1mm,left=2mm,right=2mm,top=1.5mm,bottom=1.5mm,
  title={Interpretation},fonttitle=\bfseries,#1
}
\newtcolorbox{warningbox}[1][]{
  enhanced,breakable,colback=softcoral,colframe=warmgray,
  boxrule=0.6pt,arc=1mm,left=2mm,right=2mm,top=1.5mm,bottom=1.5mm,
  title={Normalization point},fonttitle=\bfseries,#1
}

\lstdefinestyle{articlecode}{
  basicstyle=\ttfamily\small,
  backgroundcolor=\color{codeback},
  frame=single,
  rulecolor=\color{warmgray},
  breaklines=true,
  columns=fullflexible,
  keepspaces=true,
  showstringspaces=false,
  xleftmargin=1em,
  xrightmargin=1em
}
\lstset{style=articlecode}

\newcommand{\Fub}{\mathsf{F}}
\newcommand{\Fubx}{\mathsf{F}}
\newcommand{\StirlingII}[2]{\genfrac\{\}{0pt}{}{#1}{#2}}
\newcommand{\Eulerian}[2]{\genfrac{\langle}{\rangle}{0pt}{}{#1}{#2}}
\newcommand{\Coeff}[3]{[\,#1^{#2}\,]#3}
\newcommand{\Res}{\operatorname*{Res}}
\newcommand{\Li}{\operatorname{Li}}
\newcommand{\Log}{\operatorname{Log}}
\newcommand{\e}{\mathrm{e}}
\newcommand{\ii}{\mathrm{i}}
\newcommand{\dd}{\,\mathrm{d}}
\newcommand{\N}{\mathbb{N}}
\newcommand{\C}{\mathbb{C}}
\newcommand{\R}{\mathbb{R}}
\newcommand{\BellY}{\mathrm{Y}}
\newcommand{\falling}[2]{(#1)^{\underline{#2}}}
\newcommand{\GammaStar}{\Gamma^{\!*}}
\newcommand{\asympseries}{\mathrel{\sim}}
\DeclareMathOperator{\supp}{supp}

\title{\Huge\bfseries The Full Asymptotic Transseries\\[2mm]
       of the Fubini Numbers\\[3mm]
       \Large Exact Pole Sectors, General Bell-Polynomial Coefficients,\\
       Nonlinear Transforms, and Rigorous Remainders}
\author{A self-contained research article with complete derivations}
\date{September 2026}

\begin{document}
\maketitle

\begin{abstract}
The Fubini numbers, also called ordered Bell numbers or preferential-arrangement
numbers, are
\[
  \Fub_n=\sum_{j=0}^{n}j!\StirlingII{n}{j}
  \qquad(n\geq 0),
\]
and have exponential generating function $(2-\e^z)^{-1}$.  Their familiar
leading estimate
\(
 \Fub_n\sim n!/[2(\log 2)^{n+1}]
\)
uses only the nearest singularity $z=\log 2$.  The generating function actually
has the complete vertical lattice of simple poles
\(
 \zeta_k=\log 2+2\pi\ii k
\), $k\in\mathbb Z$.  Summing their contributions gives an exact, convergent
exponential transseries rather than merely a first-order singularity estimate.

Writing $\rho=\log 2$, $u_k=\rho/(\rho+2\pi\ii k)$, and
\(
 \GammaStar(n)=n!/[\sqrt{2\pi n}(n/\e)^n]
\), we prove the exact factorization
\[
 \Fub_n=
 \frac{\sqrt{2\pi n}}{2\rho}
 \left(\frac{n}{\e\rho}\right)^n
 \GammaStar(n)\sum_{k\in\mathbb Z}u_k^{n+1}
 \qquad(n\geq1).
\]
The pole sum is absolutely convergent and contains every exponentially small
oscillatory sector.  The common algebraic fluctuation factor has the all-orders
Stirling expansion
\(
 \GammaStar(n)\sim\sum_{j\ge0}c_jn^{-j}
\), where the coefficients are given explicitly by complete exponential Bell
polynomials, integer partitions, and a triangular Bernoulli recurrence.  Thus
all mixed transseries coefficients are
\(
 C_{j,k}=c_j\rho/(\rho+2\pi\ii k)
\).

We supply explicit two-parameter remainder bounds for simultaneous truncation
of the algebraic order and the pole lattice; an exact Borel--Laplace
representation of the common Stirling factor; convergent multi-index formulae
for logarithms, reciprocals, arbitrary analytic functions of the normalized
sequence, and the ratio $n\Fub_{n-1}/\Fub_n$; and a uniform extension to the
Fubini polynomials
\(
 \Fub_n(x)=\sum_j j!\StirlingII{n}{j}x^j
\), $x>0$.  A final meromorphic-EGF theorem explains why simple poles have
constant sector amplitudes and how higher-order Fubini generating functions
produce polynomial-in-$n$ amplitudes governed by generalized Bernoulli
numbers.
\end{abstract}

\begin{masterbox}
Put
\[
 \rho=\log 2,\qquad
 \zeta_k=\rho+2\pi\ii k,\qquad
 \Lambda_k=\Log\!\left(\frac{\zeta_k}{\rho}\right),
 \qquad k\in\mathbb Z,
\]
with the principal logarithm.  For every integer $n\ge1$,
\begin{align*}
 \Fub_n
 &=\frac{n!}{2}\sum_{k\in\mathbb Z}\zeta_k^{-n-1}\\
 &=\frac{n!}{2\rho^{n+1}}
   \left[1+2\sum_{k\ge1}
   \left(\frac{\rho}{|\zeta_k|}\right)^{n+1}
   \cos\!\bigl((n+1)\arg\zeta_k\bigr)\right].
\end{align*}
This is an \emph{exact convergent pole-lattice transseries}.  If
\[
 \exp\!\left(\sum_{r\ge1}
   \frac{B_{r+1}}{r(r+1)}t^r\right)
   =\sum_{j\ge0}c_jt^j,
\]
where $B_m$ are the Bernoulli numbers and the vanishing odd Bernoulli numbers
are understood, then the elementary-scale transseries is
\[
 \Fub_n\asympseries
 \frac{\sqrt{2\pi n}}{2\rho}
 \left(\frac{n}{\e\rho}\right)^n
 \sum_{k\in\mathbb Z}\e^{-n\Lambda_k}
 \sum_{j\ge0}\frac{C_{j,k}}{n^j},
 \qquad
 C_{j,k}=c_j\e^{-\Lambda_k}=c_j\frac{\rho}{\zeta_k}.
\]
The $k$-sum is convergent for every $n\ge1$; the $j$-sum is the usual
Poincar\'e expansion of the scaled gamma factor and is interpreted
asymptotically or by its positive-axis Borel sum.
\end{masterbox}

\tableofcontents

\section*{Notation and conventions}
\addcontentsline{toc}{section}{Notation and conventions}
Throughout,
\[
 \rho:=\log 2,
 \qquad \zeta_k:=\rho+2\pi\ii k\quad(k\in\mathbb Z).
\]
The principal logarithm is denoted by $\Log$.  Bernoulli numbers use
\[
 \frac{t}{\e^t-1}=\sum_{m\ge0}B_m\frac{t^m}{m!},
 \qquad B_1=-\frac12.
\]
The complete exponential Bell polynomials $\BellY_j$ are normalized by
\[
 \exp\!\left(\sum_{r\ge1}x_r\frac{t^r}{r!}\right)
 =\sum_{j\ge0}\BellY_j(x_1,\ldots,x_j)\frac{t^j}{j!}.
\]
A sum over $k\in\mathbb Z$ is ordinary and absolutely convergent whenever the
summand is $O(|k|^{-1-\varepsilon})$.  At exponent one we write
$\sum_{k\in\mathbb Z}^{\mathrm{sym}}$ for the symmetric limit
$\lim_{K\to\infty}\sum_{k=-K}^{K}$.

The symbol $\sim$ has two related uses.  Between functions it denotes
asymptotic equivalence.  Between a function and a formal inverse-power series
it denotes a Poincar\'e asymptotic expansion.  All exact identities are marked
with equality, in particular the complete pole-lattice sum.

\section{Introduction and scope}\label{sec:introduction}

\subsection{From one dominant pole to a complete transseries}
The elementary derivation of the leading asymptotic formula for $\Fub_n$ is a
standard application of singularity analysis.  The nearest singularity of
$(2-\e^z)^{-1}$ to the origin is the positive real pole $\rho=\log2$, whose
residue is $-1/2$.  Therefore
\[
 \Fub_n\sim\frac{n!}{2\rho^{n+1}}.
\]
This statement is correct but discards far more structure than it retains.
Every solution of $\e^z=2$ is a pole, and those solutions form the lattice
\[
 \rho+2\pi\ii\mathbb Z.
\]
The nonreal poles occur in conjugate pairs.  Their coefficient contributions
are exponentially smaller than the real-pole contribution and oscillate with
$n$.  If $q_k:=\rho/|\rho+2\pi\ii k|$, the principal-pole approximation
has error $O(q_1^{n+1})$, whereas retaining the first conjugate pair leaves
$O(q_2^{n+1})$; the improvement factor is therefore
$O((q_2/q_1)^{n+1})$.  Retaining the entire lattice gives an exact formula.

The word \emph{full} in the title is used in the following precise sense.
\begin{enumerate}
 \item Every singularity of the Fubini exponential generating function is
       represented by a transseries sector.
 \item Every algebraic inverse-power coefficient multiplying every sector is
       given in closed form and by an effective recurrence.
 \item Fixed finite truncations in both directions come with rigorous error
       estimates.
 \item The factorial fluctuation may either be retained exactly, expanded to
       arbitrary algebraic order, or reconstructed by an exact Borel--Laplace
       integral on the positive axis.
\end{enumerate}
Thus no pole sector is omitted.  The only divergent object is the familiar
Stirling series that appears when the exact factor $n!$ is replaced by an
elementary expression.  Keeping $n!$ or $\GammaStar(n)$ exact removes even
that divergence.

\subsection{Three normalizations that should not be conflated}
It is useful to separate three layers:
\begin{align*}
 \text{exact factorial normalization:}\quad
 &\frac{2\rho^{n+1}}{n!}\Fub_n,\\
 \text{elementary normalization:}\quad
 &\frac{2\rho}{\sqrt{2\pi n}}
   \left(\frac{\e\rho}{n}\right)^n\Fub_n,\\
 \text{logarithmic normalization:}\quad
 &\log\Fub_n-\bigl((n+\tfrac12)\log n-n-(n+1)\log\rho
      +\tfrac12\log(2\pi)-\log2\bigr).
\end{align*}
In the first normalization the answer is an exact convergent sum of
exponentials with constant amplitudes.  In the second, every exponential
sector is multiplied by the same Stirling series.  In the third, the
algebraic and exponential corrections add, but products of pole sectors
appear because of the logarithm.  A substantial part of the paper is devoted
to making these translations explicit.

\begin{warningbox}
A simple pole contributes \emph{one exact constant-amplitude sector}.  It does
not produce an independent infinite $1/n$ fluctuation series.  The inverse
powers in the elementary normalization come solely from replacing $n!$ by
Stirling's expansion.  Higher-order poles do produce polynomial-in-$n$
sector amplitudes; this distinction is proved in \cref{sec:meromorphic}.
\end{warningbox}

\subsection{Relation to coefficient calculus}
The derivation uses a small number of reusable coefficient-calculus
principles:
\begin{enumerate}
 \item labelled set partitions turn the Stirling sum into a geometric series
       of exponential generating functions;
 \item the Mittag--Leffler decomposition turns a meromorphic generating
       function into a sum of elementary coefficient contributions;
 \item exponentiating the Bernoulli expansion of $\log\GammaStar$ turns
       logarithmic coefficients into complete Bell polynomials;
 \item applying an analytic function to a convergent exponential correction
       turns ordinary powers into multinomial or multi-index coefficients.
\end{enumerate}
These are exactly the points at which Stirling numbers, Bell polynomials,
formal composition, and transseries algebra meet.

\section{Combinatorial and probabilistic foundations}\label{sec:foundations}

\subsection{Ordered set partitions}

\begin{definition}
An \emph{ordered set partition} of an $n$-element labelled set is a sequence
$(B_1,\ldots,B_j)$ of nonempty pairwise disjoint blocks whose union is the
whole set.  Equivalently, it is a weak ordering: elements in the same block
are tied, while the blocks themselves are linearly ordered.
\end{definition}

\begin{theorem}[Stirling transform]\label{thm:stirling-sum}
The number of ordered set partitions of an $n$-element labelled set is
\[
 \Fub_n=\sum_{j=0}^{n}j!\StirlingII{n}{j}.
\]
In particular,
\[
 \Fub_0=1,\quad \Fub_1=1,\quad \Fub_2=3,\quad \Fub_3=13,
 \quad \Fub_4=75,
\]
and the sequence continues $541,4683,47293,545835,\ldots$.
\end{theorem}

\begin{proof}
First partition the labelled set into $j$ nonempty unordered blocks.  There
are $\StirlingII{n}{j}$ choices.  Then put the $j$ blocks in one of their $j!$
linear orders.  Summation over every possible number of blocks gives the
formula.  The empty set has one empty ordered partition, accounting for
$\Fub_0=1$.
\end{proof}

\subsection{Exponential generating function}

\begin{theorem}[Fubini EGF]\label{thm:egf}
The exponential generating function is
\[
 \mathcal F(z):=\sum_{n\ge0}\Fub_n\frac{z^n}{n!}
 =\frac{1}{2-\e^z}.
\]
\end{theorem}

\begin{proof}
For fixed $j$, the standard exponential generating function of partitions
into $j$ nonempty blocks is
\[
 \sum_{n\ge0}\StirlingII{n}{j}\frac{z^n}{n!}
 =\frac{(\e^z-1)^j}{j!}.
\]
Multiplication by $j!$ orders the blocks.  Therefore, as a formal power series,
\[
 \mathcal F(z)
 =\sum_{j\ge0}(\e^z-1)^j
 =\frac{1}{1-(\e^z-1)}
 =\frac{1}{2-\e^z}.
\]
The same calculation is analytic whenever $|\e^z-1|<1$, and analytic
continuation gives the displayed meromorphic function.
\end{proof}

\begin{corollary}[Binomial recurrence]\label{cor:recurrence}
For $n\ge1$,
\[
 \Fub_n=\sum_{j=0}^{n-1}\binom{n}{j}\Fub_j
       =\sum_{j=1}^{n}\binom{n}{j}\Fub_{n-j}.
\]
\end{corollary}

\begin{proof}
The identity $(2-\e^z)\mathcal F(z)=1$ gives, in positive degree,
\[
 2\Fub_n-\sum_{j=0}^{n}\binom{n}{j}\Fub_j=0.
\]
The $j=n$ summand is $\Fub_n$; moving it to the other side proves the first
form.  Replacing $j$ by $n-j$ gives the second.
\end{proof}

\subsection{A Dobinski-type moment formula}
The geometric expansion of the EGF already anticipates the pole formula:
\[
 \frac{1}{2-\e^z}
 =\frac12\frac{1}{1-\e^z/2}
 =\frac12\sum_{m\ge0}\frac{\e^{mz}}{2^m}.
\]
Coefficient extraction gives the following exact analogue of Dobinski's
formula for ordinary Bell numbers.

\begin{theorem}[Geometric moment representation]\label{thm:moment}
For every $n\ge0$,
\[
 \Fub_n=\frac12\sum_{m=0}^{\infty}\frac{m^n}{2^m},
\]
where $0^0=1$.  For $n\ge1$ this is equivalently
\[
 \Fub_n=\frac12\Li_{-n}\!\left(\frac12\right).
\]
\end{theorem}

\begin{proof}
Expand each exponential in the preceding geometric series and interchange the
absolutely convergent sums in a neighborhood of $z=0$:
\[
 \frac12\sum_{m\ge0}2^{-m}\e^{mz}
 =\sum_{n\ge0}\left(\frac12\sum_{m\ge0}\frac{m^n}{2^m}\right)
  \frac{z^n}{n!}.
\]
Comparison with \cref{thm:egf} proves the result.  When $n\ge1$ the $m=0$
term vanishes, and the remaining series is the defining series for the
polylogarithm at negative integer order.
\end{proof}

\begin{interpretbox}
Let $G$ be the geometric random variable
\[
 \Pr(G=m)=2^{-m-1},\qquad m=0,1,2,\ldots.
\]
Then $\Fub_n=\mathbb E[G^n]$.  The Fubini numbers are therefore a genuine Stieltjes
moment sequence, not merely a combinatorial counting sequence.  The same
probability model will yield the weighted Fubini polynomials in
\cref{sec:weighted}.
\end{interpretbox}

\subsection{Eulerian form and cumulant Bell polynomials}
Let $\Eulerian{n}{j}$ be the Eulerian number counting
permutations of $n$ elements with $j$ descents, and put
\[
 A_n(t)=\sum_{j=0}^{n-1}\Eulerian{n}{j} t^j
 \qquad(n\ge1).
\]
The classical rational form of the negative polylogarithm is
\[
 \Li_{-n}(t)=\frac{tA_n(t)}{(1-t)^{n+1}}.
\]
Using the symmetry of the Eulerian numbers gives
\begin{equation}\label{eq:eulerian-form}
 \Fub_n=2^{n-1}A_n\!\left(\frac12\right)
 =\sum_{j=0}^{n-1}2^j\Eulerian{n}{j},
 \qquad n\ge1.
\end{equation}

The moment generating function of $G$ is $\mathcal F(z)$.  Its cumulant
generating function is
\[
 K(z)=\log\mathcal F(z)
 =-\log2-\log\!\left(1-\frac{\e^z}{2}\right)
 =\sum_{m\ge1}\frac{\e^{mz}}{m2^m}-\log2.
\]
Consequently the $r$th cumulant is
\begin{equation}\label{eq:cumulants}
 \kappa_r=K^{(r)}(0)=\Li_{1-r}\!\left(\frac12\right),
 \qquad r\ge1.
\end{equation}
Moment--cumulant conversion now gives a second Bell-polynomial formula.

\begin{proposition}[Moment--cumulant formula]\label{prop:moment-cumulant}
For $n\ge0$,
\[
 \Fub_n=\BellY_n(\kappa_1,\ldots,\kappa_n),
 \qquad
 \kappa_r=\Li_{1-r}\!\left(\frac12\right).
\]
For $r\ge2$, $\kappa_r=2\Fub_{r-1}$, while $\kappa_1=\Fub_0=1$.
\end{proposition}

\begin{proof}
The defining identity for complete exponential Bell polynomials is precisely
\[
 \exp\!\left(\sum_{r\ge1}\kappa_r\frac{z^r}{r!}\right)
 =\sum_{n\ge0}\BellY_n(\kappa_1,\ldots,\kappa_n)\frac{z^n}{n!}.
\]
The left-hand side is $\exp K(z)=\mathcal F(z)$, so coefficient comparison
proves the first formula.  The relation with $\Fub_{r-1}$ follows from
\cref{thm:moment}, noting separately the exceptional $0^0$ contribution at
$r=1$.
\end{proof}


\section{The exact pole lattice}\label{sec:poles}

\subsection{Location and residues}
The singularities of $\mathcal F(z)=(2-\e^z)^{-1}$ are the solutions of
$\e^z=2$:
\[
 \zeta_k=\rho+2\pi\ii k,\qquad k\in\mathbb Z.
\]
Since
\[
 \frac{\dd}{\dd z}(2-\e^z)\bigg|_{z=\zeta_k}=-2,
\]
every pole is simple and
\begin{equation}\label{eq:residue}
 \Res_{z=\zeta_k}\mathcal F(z)=-\frac12.
\end{equation}
The moduli $|\zeta_k|$ increase strictly with $|k|$.  Thus the real pole
$\zeta_0=\rho$ is dominant, followed by the conjugate pair
$\zeta_{\pm1}$, then $\zeta_{\pm2}$, and so forth.

\begin{figure}[htbp]
\centering
\begin{tikzpicture}[scale=0.72,>=Latex]
  \draw[->,warmgray] (-1.1,0)--(4.7,0) node[right] {$\Re z$};
  \draw[->,warmgray] (0,-4.5)--(0,4.7) node[above] {$\Im z$};
  \draw[dashed,darkteal] (1.25,-4.2)--(1.25,4.2);
  \node[below,darkteal] at (1.25,0) {$\rho$};
  \foreach \j in {-3,-2,-1,0,1,2,3}{
    \fill[deepolive] (1.25,{1.25*\j}) circle (2.1pt);
  }
  \node[right] at (1.25,0.12) {$\zeta_0$};
  \node[right] at (1.25,1.25) {$\zeta_1$};
  \node[right] at (1.25,-1.25) {$\zeta_{-1}$};
  \draw[->,linkblue,thick] (0,0)--(1.25,1.25);
  \draw[linkblue] (0.53,0) arc[start angle=0,end angle=45,radius=0.53];
  \node[linkblue] at (0.72,0.25) {$\theta_1$};
  \node[linkblue,above left] at (0.7,0.75) {$R_1$};
  \draw[decorate,decoration={brace,amplitude=4pt},warmgray]
    (1.55,1.25)--(1.55,2.5) node[midway,right=5pt] {$2\pi$};
\end{tikzpicture}
\caption{The poles of $(2-\e^z)^{-1}$ lie on the vertical line
$\Re z=\rho$.  Their equal residues and increasing moduli make the pole
ordering identical to the exponential-sector ordering.  The drawing rescales
the vertical direction for legibility.}
\label{fig:pole-lattice}
\end{figure}

\subsection{A global Mittag--Leffler decomposition}
Put $w=z-\rho$.  Since $\e^z=2\e^w$,
\begin{equation}\label{eq:coth-start}
 \mathcal F(z)=\frac{1}{2(1-\e^w)}
 =\frac14-\frac14\coth\!\left(\frac w2\right).
\end{equation}
The paired partial-fraction expansion
\begin{equation}\label{eq:coth-ml}
 \coth\!\left(\frac w2\right)
 =\frac{2}{w}+4w\sum_{k\ge1}\frac{1}{w^2+4\pi^2k^2}
 =2\sum_{k\in\mathbb Z}^{\mathrm{sym}}\frac{1}{w-2\pi\ii k}
\end{equation}
is locally uniformly convergent away from its poles.  Combining
\cref{eq:coth-start,eq:coth-ml} yields the exact global decomposition.

\begin{theorem}[Exact Mittag--Leffler expansion]\label{thm:mittag}
For $z\notin\rho+2\pi\ii\mathbb Z$,
\begin{equation}\label{eq:mittag}
 \frac{1}{2-\e^z}
 =\frac14-\frac12\sum_{k\in\mathbb Z}^{\mathrm{sym}}
   \frac{1}{z-\zeta_k}.
\end{equation}
The paired series converges locally uniformly on the complement of the pole
lattice.
\end{theorem}

\begin{proof}
Pairing $k$ with $-k$ gives
\[
 \frac{1}{w-2\pi\ii k}+\frac{1}{w+2\pi\ii k}
 =\frac{2w}{w^2+4\pi^2k^2}=O(k^{-2})
\]
uniformly on compact $w$-sets.  Hence the paired sum is locally normally
convergent away from the poles.  The classical product
\(
 \sinh(w/2)=(w/2)\prod_{k\ge1}(1+w^2/(4\pi^2k^2))
\)
and logarithmic differentiation give \cref{eq:coth-ml}.  Substitution into
\cref{eq:coth-start} proves \cref{eq:mittag}.
\end{proof}

\subsection{Exact coefficient extraction}
For a fixed nonzero $\zeta$,
\[
 -\frac{1}{2(z-\zeta)}
 =\frac{1}{2\zeta}\frac{1}{1-z/\zeta}
 =\frac12\sum_{n\ge0}\frac{z^n}{\zeta^{n+1}}.
\]
When $n\ge1$, the resulting sum over $k$ is absolutely convergent because its
terms are $O(|k|^{-n-1})$.

\begin{theorem}[Exact all-pole coefficient formula]\label{thm:exact-poles}
For every integer $n\ge1$,
\begin{equation}\label{eq:exact-poles}
 \boxed{\displaystyle
 \Fub_n=\frac{n!}{2}\sum_{k\in\mathbb Z}
       \frac{1}{(\rho+2\pi\ii k)^{n+1}}.}
\end{equation}
For $n=0$, the regularized unified formula is
\begin{equation}\label{eq:nzero-regularized}
 \Fub_0=\frac14+\frac12\sum_{k\in\mathbb Z}^{\mathrm{sym}}
       \frac{1}{\rho+2\pi\ii k}=1.
\end{equation}
\end{theorem}

\begin{proof}
Take the coefficient of $z^n$ in \cref{eq:mittag} and multiply by $n!$.
For $n\ge1$, absolute convergence justifies termwise coefficient extraction.
For $n=0$, the symmetric sum is
\[
 \sum_{k\in\mathbb Z}^{\mathrm{sym}}\frac{1}{\rho+2\pi\ii k}
 =\frac12\coth\!\left(\frac\rho2\right).
\]
Because $\e^\rho=2$, one has
\(
 \coth(\rho/2)=(2+1)/(2-1)=3
\), so the right-hand side of \cref{eq:nzero-regularized} is
$1/4+3/4=1$.
\end{proof}

\begin{remark}
The separate constant $1/4$ matters only at $n=0$.  This is why the clean
bilateral pole formula is usually stated for $n\ge1$.  No regularization is
needed at any positive order.
\end{remark}

\subsection{A second proof by Poisson summation}
The geometric-moment formula supplies an independent route to the same
identity and simultaneously permits a continuous order parameter.

\begin{theorem}[Lipschitz--Poisson formula]\label{thm:poisson}
For $\nu\in\mathbb C$ with $\Re\nu>0$, define
\[
 \mathscr F(\nu):=\frac12\sum_{m\ge1}m^\nu\e^{-\rho m}
 =\frac12\Li_{-\nu}\!\left(\frac12\right),
\]
using the real logarithm in $m^\nu$.  Then
\begin{equation}\label{eq:continuous-poles}
 \mathscr F(\nu)
 =\frac{\Gamma(\nu+1)}{2}
  \sum_{k\in\mathbb Z}(\rho+2\pi\ii k)^{-\nu-1},
\end{equation}
where the complex powers use the principal logarithm.  At every integer
$\nu=n\ge1$, this is \cref{eq:exact-poles}.
\end{theorem}

\begin{proof}
Set
\[
 f_\nu(x)=\begin{cases}
 x^\nu\e^{-\rho x},&x>0,\\
 0,&x\le0.
 \end{cases}
\]
For $\Re\nu>0$ the function is continuous at the origin, integrable, and has
sufficient decay and local regularity for the classical Poisson summation
formula (or for its standard Abel-regularized form followed by a limit).  With
the Fourier-transform convention
\(
 \widehat f(\xi)=\int_{\mathbb R}f(x)\e^{-2\pi\ii x\xi}\dd x
\),
\[
 \widehat f_\nu(\xi)
 =\int_0^\infty x^\nu\e^{-(\rho+2\pi\ii\xi)x}\dd x
 =\frac{\Gamma(\nu+1)}{(\rho+2\pi\ii\xi)^{\nu+1}}.
\]
The Fourier series is absolutely convergent because
$\Re(\nu+1)>1$.  Hence
\[
 \sum_{m\in\mathbb Z}f_\nu(m)
 =\sum_{k\in\mathbb Z}\widehat f_\nu(k).
\]
The left side is $\sum_{m\ge1}m^\nu\e^{-\rho m}$.  Division by two proves
\cref{eq:continuous-poles}.
\end{proof}

\begin{interpretbox}
The exact pole transseries is not an accident restricted to integer
coefficients.  It is the sampling duality between an exponentially damped
power on the positive half-line and a vertical reciprocal-power lattice in
the Fourier domain.  The integer Fubini numbers are special values of this
continuous-order identity.
\end{interpretbox}

\subsection{Hurwitz-zeta and differential forms}
Let
\[
 a:=\frac{\rho}{2\pi\ii}=-\frac{\ii\rho}{2\pi}.
\]
Splitting the bilateral sum into $k\ge0$ and $k\le-1$ gives, for $n\ge1$,
\begin{equation}\label{eq:hurwitz-form}
 \Fub_n=\frac{n!}{2(2\pi\ii)^{n+1}}
 \left[
  \zeta(n+1,a)+(-1)^{n+1}\zeta(n+1,1-a)
 \right],
\end{equation}
where the zeta symbols on the right are Hurwitz zeta functions.  This form is
occasionally convenient for symbolic manipulation, although the conjugate
pole-pair form developed below is numerically more transparent.

Define
\[
 S_s(\rho):=\sum_{k\in\mathbb Z}^{\mathrm{sym}}
             (\rho+2\pi\ii k)^{-s}.
\]
Then
\[
 S_1(\rho)=\frac12\coth\!\left(\frac\rho2\right),
 \qquad
 S_{s+1}(\rho)=-\frac1s\frac{\dd}{\dd\rho}S_s(\rho).
\]
Consequently, for $n\ge1$,
\begin{equation}\label{eq:coth-derivative}
 \Fub_n=\frac{(-1)^n}{4}
 \left.\frac{\dd^n}{\dd x^n}\coth\!\left(\frac x2\right)
 \right|_{x=\rho}.
\end{equation}
The derivative-polynomial description of $\coth$ is another route from the
pole lattice to Eulerian polynomials and hence to \cref{eq:eulerian-form}.

\section{The complete exponential transseries}\label{sec:exp-transseries}

\subsection{Complex and real sector coordinates}
Define
\begin{equation}\label{eq:u-def}
 u_k:=\frac{\rho}{\zeta_k}
 =\frac{\rho}{\rho+2\pi\ii k},
 \qquad k\in\mathbb Z.
\end{equation}
Thus $u_0=1$, $u_{-k}=\overline{u_k}$, and $|u_k|<1$ for $k\ne0$.
For $k\ge1$, write
\begin{align}
 R_k&:=|\zeta_k|=\sqrt{\rho^2+4\pi^2k^2},\label{eq:Rk}\\
 q_k&:=|u_k|=\frac{\rho}{R_k},\label{eq:qk}\\
 \alpha_k&:=-\log q_k
 =\frac12\log\!\left(1+\frac{4\pi^2k^2}{\rho^2}\right),\label{eq:alphak}\\
 \theta_k&:=\arg\zeta_k=\arctan\!\left(\frac{2\pi k}{\rho}\right).
 \label{eq:thetak}
\end{align}
Then
\[
 u_k=q_k\e^{-\ii\theta_k}
     =\e^{-\alpha_k-\ii\theta_k},
 \qquad
 u_{-k}=q_k\e^{\ii\theta_k}.
\]
The real actions satisfy
\(
 0<\alpha_1<\alpha_2<\cdots
\) and $\alpha_k\to\infty$.

\begin{theorem}[Exact pole-lattice transseries]\label{thm:exact-transseries}
For every $n\ge1$, put
\[
 T_n:=\frac{2\rho^{n+1}}{n!}\Fub_n.
\]
Then the following are exact identities:
\begin{align}
 T_n
 &=\sum_{k\in\mathbb Z}u_k^{n+1},\label{eq:T-complex}\\
 &=1+2\sum_{k\ge1}q_k^{n+1}
       \cos\!\bigl((n+1)\theta_k\bigr),\label{eq:T-real}\\
 \Fub_n
 &=\frac{n!}{2\rho^{n+1}}
   \left[1+2\sum_{k\ge1}\e^{-(n+1)\alpha_k}
       \cos\!\bigl((n+1)\theta_k\bigr)\right].
 \label{eq:exact-real-transseries}
\end{align}
Both series are absolutely convergent.
\end{theorem}

\begin{proof}
Multiply \cref{eq:exact-poles} by $2\rho^{n+1}/n!$ to obtain
\cref{eq:T-complex}.  Pairing $k$ and $-k$ gives
\[
 u_k^{n+1}+u_{-k}^{n+1}
 =2q_k^{n+1}\cos((n+1)\theta_k),
\]
which proves the real forms.  Absolute convergence follows from
$q_k^{n+1}=O(k^{-n-1})$.
\end{proof}

\begin{remark}[Why this is a transseries]
Each $k\ne0$ contributes a transmonomial
\(
 \e^{-n\alpha_{|k|}}\e^{\mp\ii n\theta_{|k|}}
\)
with an additional constant factor $u_k$.  The support is locally finite in
real action: below any fixed action bound there are only finitely many $k$.
The conjugate-pair form is real, while the complex form is algebraically more
convenient.  Unlike a generic formal transseries, this one converges and sums
exactly to the normalized integer sequence.
\end{remark}

\subsection{The first sectors}
The first six action parameters are shown in \cref{tab:actions}.  Angles are in
radians.

\begin{table}[htbp]
\centering
\caption{Pole moduli, sector ratios, real actions, and phases for
$\rho=\log2$.}
\label{tab:actions}
\begin{tabular}{@{}rllll@{}}
\toprule
$k$ & $R_k$ & $q_k$ & $\alpha_k$ & $\theta_k$\\
\midrule
1 & 6.32130292093929 & 0.109652580999385 & 2.21043826586918 & 1.46092281020826\\
2 & 12.5854727138613 & 0.055075180433750 & 2.89905611015294 & 1.51569326525480\\
3 & 18.8622960281387 & 0.036747762813494 & 3.30367793099960 & 1.53404028826676\\
4 & 25.1422977208456 & 0.027568967174597 & 3.59106451657788 & 1.54322386613582\\
5 & 31.4235722527350 & 0.022058191697146 & 3.81407124054088 & 1.54873634591964\\
6 & 37.7054835106352 & 0.018383193000679 & 3.99631845566911 & 1.55241209822855\\
\bottomrule
\end{tabular}
\end{table}

Thus the beginning of the exact expansion is
\begin{align}\label{eq:first-sectors}
 \Fub_n=\frac{n!}{2\rho^{n+1}}\bigl[&1
 +2q_1^{n+1}\cos((n+1)\theta_1)
 +2q_2^{n+1}\cos((n+1)\theta_2)\nonumber\\
 &+2q_3^{n+1}\cos((n+1)\theta_3)+\cdots\bigr].
\end{align}
The leading correction is not of fixed sign.  Since $\theta_1$ is close to
$\pi/2$, its sign exhibits a near four-step oscillation with a slowly drifting
phase.  This is why pointwise error ratios can be much smaller than their
envelopes at accidental cancellations.

\subsection{Rigorous truncation bounds}
For $K\ge0$, define the $K$-pair truncation
\begin{equation}\label{eq:TK}
 T_{n,K}:=1+2\sum_{k=1}^{K}q_k^{n+1}
                   \cos((n+1)\theta_k),
\end{equation}
where an empty sum means $T_{n,0}=1$.

\begin{theorem}[Pole-tail bounds]\label{thm:pole-tail}
For $n\ge1$ and $K\ge0$,
\begin{align}
 |T_n-T_{n,K}|
 &\le B_{n,K}:=2\sum_{k>K}q_k^{n+1},\label{eq:B-exact}\\
 &\le 2\left(\frac{\rho}{2\pi}\right)^{n+1}
       \zeta(n+1,K+1),\label{eq:B-hurwitz}
\end{align}
where $\zeta(s,a)$ is the Hurwitz zeta function.  If $K\ge1$, then also
\begin{equation}\label{eq:B-integral}
 |T_n-T_{n,K}|
 \le \frac{2}{n}\left(\frac{\rho}{2\pi}\right)^{n+1}K^{-n}.
\end{equation}
For the leading approximation $K=0$ one has the uniform estimate
\begin{equation}\label{eq:global-relative-bound}
 |T_n-1|\le\frac{\rho^2}{12}=0.0400377511598501\ldots
 \qquad(n\ge1).
\end{equation}
\end{theorem}

\begin{proof}
The triangle inequality applied to the omitted conjugate pairs gives
\cref{eq:B-exact}.  Since
\[
 q_k=\frac{\rho}{\sqrt{\rho^2+4\pi^2k^2}}
 \le\frac{\rho}{2\pi k},
\]
summing the majorant proves \cref{eq:B-hurwitz}.  For $K\ge1$ and $s=n+1$,
\[
 \sum_{k=K+1}^{\infty}k^{-s}
 \le\int_K^\infty x^{-s}\dd x=\frac{K^{-n}}{n},
\]
which proves \cref{eq:B-integral}.  Finally,
$\rho/(2\pi)<1$ and $\zeta(n+1)\le\zeta(2)$, so
\[
 B_{n,0}
 \le2\left(\frac{\rho}{2\pi}\right)^2\zeta(2)
 =2\frac{\rho^2}{4\pi^2}\frac{\pi^2}{6}
 =\frac{\rho^2}{12}.
\]
\end{proof}

\begin{corollary}[Successive asymptotic levels]\label{cor:successive}
As $n\to\infty$,
\begin{align}
 \Fub_n
 &\sim\frac{n!}{2\rho^{n+1}},\label{eq:leading-asymp}\\
 \Fub_n
 &=\frac{n!}{2\rho^{n+1}}
 \left[1+2q_1^{n+1}\cos((n+1)\theta_1)
 +O(q_2^{n+1})\right].\label{eq:first-correction}
\end{align}
More generally, after $K$ conjugate pairs the relative remainder is bounded by
\cref{eq:B-hurwitz} and is therefore
$O((\rho/(2\pi(K+1)))^{n+1})$ for fixed $K$.
\end{corollary}

\begin{proof}
The first statement follows from \cref{eq:global-relative-bound} together with
$B_{n,0}\to0$.  For the second, separate the $k=1$ pair and apply
\cref{eq:B-hurwitz} with $K=1$.  The general case is identical.
\end{proof}

\subsection{No hidden local inverse-power coefficients}
Near a pole $\zeta_k$,
\[
 \mathcal F(z)=-\frac{1}{2(z-\zeta_k)}+H_k(z),
\]
where $H_k$ is analytic near $\zeta_k$.  The singular term contributes exactly
\(
 \tfrac12\zeta_k^{-n-1}
\)
to the ordinary Taylor coefficient.  There is no expansion of the residue in
$n^{-1}$.  Once the singularity at $\zeta_k$ has been subtracted, the analytic
remainder is controlled by other singularities, not by an infinite local
series at the same pole.  This elementary observation explains the constant
sector amplitudes in \cref{eq:T-complex}.  The inverse-power coefficients that
appear after elementary normalization will come entirely from $n!$.


\section{The factorial fluctuation and its general coefficients}
\label{sec:stirling-layer}

\subsection{Exact elementary normalization}
Define the scaled gamma function for $z>0$ by
\begin{equation}\label{eq:scaled-gamma-def}
 \GammaStar(z):=
 \frac{\Gamma(z)}{\sqrt{2\pi}\,z^{z-1/2}\e^{-z}}.
\end{equation}
At a positive integer $n$ this may be written in the especially useful form
\begin{equation}\label{eq:factorial-scaled}
 n!=\sqrt{2\pi n}\left(\frac n\e\right)^n\GammaStar(n).
\end{equation}
Introduce the elementary leading scale
\begin{equation}\label{eq:A-def}
 \mathcal A(n):=
 \frac{\sqrt{2\pi n}}{2\rho}
 \left(\frac{n}{\e\rho}\right)^n.
\end{equation}
Combining \cref{eq:factorial-scaled,eq:T-complex} gives the exact factorization
\begin{equation}\label{eq:exact-factorization}
 \boxed{\displaystyle
 \Fub_n=\mathcal A(n)\GammaStar(n)T_n,
 \qquad
 T_n=\sum_{k\in\mathbb Z}u_k^{n+1}.}
\end{equation}
The transseries problem has therefore split into two independent pieces:
$T_n$ is an exact convergent exponential sum, while $\GammaStar(n)$ carries the
common inverse-power fluctuation.

\subsection{Binet's exact Borel--Laplace representation}
For $\Re z>0$, Binet's formula can be written as
\begin{equation}\label{eq:binet-laplace}
 \log\GammaStar(z)
 =\int_0^\infty \e^{-zt}h(t)\dd t,
 \qquad
 h(t):=\frac1t\left(\frac{1}{\e^t-1}-\frac1t+\frac12\right).
\end{equation}
The apparent singularity of $h$ at $t=0$ is removable.  Its Taylor expansion is
\begin{equation}\label{eq:h-series}
 h(t)=\sum_{m\ge1}\frac{B_{2m}}{(2m)!}t^{2m-2}
 =\frac1{12}-\frac{t^2}{720}+\frac{t^4}{30240}-\cdots.
\end{equation}
The nearest complex singularities of $h$ are at $t=\pm2\pi\ii$.

\begin{proposition}[Exact summation of the logarithmic Stirling series]
\label{prop:binet-borel}
For $z>0$, \cref{eq:binet-laplace} is an exact positive-axis
Borel--Laplace representation.  Watson's lemma gives
\begin{equation}\label{eq:log-gammastar-asymp}
 \log\GammaStar(z)
 \asympseries
 \sum_{m\ge1}\frac{B_{2m}}{2m(2m-1)z^{2m-1}}
 =\sum_{r\ge1}\frac{b_r}{z^r},
\end{equation}
where
\begin{equation}\label{eq:b-def}
 b_r:=\frac{B_{r+1}}{r(r+1)}.
\end{equation}
Thus $b_{2m}=0$ and
$\displaystyle b_{2m-1}=B_{2m}/[2m(2m-1)]$.
\end{proposition}

\begin{proof}
The expansion \cref{eq:h-series} follows from the Bernoulli generating
function after the singular terms $t^{-1}-1/2$ are removed.  Termwise Laplace
integration formally gives
\[
 \int_0^\infty \e^{-zt}
 \frac{B_{2m}}{(2m)!}t^{2m-2}\dd t
 =\frac{B_{2m}(2m-2)!}{(2m)!z^{2m-1}}
 =\frac{B_{2m}}{2m(2m-1)z^{2m-1}}.
\]
Watson's lemma makes this a Poincar\'e expansion.  The exact integral is
analytic and unambiguous on the positive Laplace ray because the poles of
$h$ lie off that ray.  Binet's formula identifies it with
$\log\GammaStar(z)$.
\end{proof}

\begin{remark}
The singularities $\pm2\pi\ii,\pm4\pi\ii,\ldots$ of the Borel kernel control
the factorial divergence and hyperasymptotic behavior of Stirling's series.
They are conceptually distinct from the poles
$\rho+2\pi\ii k$ of the Fubini EGF.  Formula \cref{eq:exact-factorization}
keeps these two singularity geometries separate.
\end{remark}

\subsection{Complete Bell-polynomial formula for the Stirling coefficients}
Define $c_j$ by the formal exponential
\begin{equation}\label{eq:c-generating}
 C(t):=\exp\!\left(\sum_{r\ge1}b_rt^r\right)
 =\sum_{j\ge0}c_jt^j.
\end{equation}
Then
\begin{equation}\label{eq:gammastar-series}
 \GammaStar(n)\asympseries\sum_{j\ge0}\frac{c_j}{n^j}.
\end{equation}
The next theorem gives several equivalent all-orders formulae.

\begin{theorem}[General Stirling coefficient formulae]\label{thm:cj-formulas}
Let $c_0=1$.  For every $j\ge1$,
\begin{align}
 c_j
 &=\Coeff{t}{j}{\exp\!\left(\sum_{r\ge1}
        \frac{B_{r+1}}{r(r+1)}t^r\right)},
 \label{eq:c-coeff}\\
 &=\frac1{j!}\BellY_j\bigl(1!b_1,2!b_2,\ldots,j!b_j\bigr),
 \label{eq:c-bell}\\
 &=\sum_{\substack{m_1,\ldots,m_j\ge0\\
                    m_1+2m_2+\cdots+jm_j=j}}
     \prod_{r=1}^{j}\frac{b_r^{m_r}}{m_r!},
 \label{eq:c-partitions}\\
 &=\sum_{\substack{\nu_1,\nu_2,\ldots\ge0\\
       \sum_{m\ge1}(2m-1)\nu_m=j}}
   \prod_{m\ge1}\frac{1}{\nu_m!}
   \left(\frac{B_{2m}}{2m(2m-1)}\right)^{\nu_m}.
 \label{eq:c-odd-partitions}
\end{align}
Only finitely many factors occur in \cref{eq:c-odd-partitions}.  The
coefficients also satisfy the triangular recurrences
\begin{align}
 c_j&=\frac1j\sum_{r=1}^{j}
       \frac{B_{r+1}}{r+1}c_{j-r},
 \label{eq:c-recurrence}\\
 c_j&=\frac1j\sum_{m=1}^{\lfloor(j+1)/2\rfloor}
       \frac{B_{2m}}{2m}c_{j-2m+1}.
 \label{eq:c-recurrence-odd}
\end{align}
In particular every $c_j$ is rational and is effectively computable from
$c_0,\ldots,c_{j-1}$.
\end{theorem}

\begin{proof}
Formula \cref{eq:c-coeff} is the definition \cref{eq:c-generating} with
\cref{eq:b-def}.  Comparing \cref{eq:c-generating} with the defining
generating function of the complete exponential Bell polynomial gives
\cref{eq:c-bell}.  Expanding each factor
\(
 \exp(b_rt^r)=\sum_{m_r\ge0}b_r^{m_r}t^{rm_r}/m_r!
\)
and collecting the coefficient of $t^j$ gives
\cref{eq:c-partitions}.  Since $b_r=0$ for even $r$, reindexing odd parts as
$r=2m-1$ gives \cref{eq:c-odd-partitions}.

For the recurrence, logarithmically differentiate \cref{eq:c-generating}:
\[
 C'(t)=C(t)\sum_{r\ge1}r b_rt^{r-1}.
\]
The coefficient of $t^{j-1}$ is
\[
 jc_j=\sum_{r=1}^{j}r b_r c_{j-r}
 =\sum_{r=1}^{j}\frac{B_{r+1}}{r+1}c_{j-r},
\]
which proves \cref{eq:c-recurrence}.  Deleting the zero even-$r$ terms gives
\cref{eq:c-recurrence-odd}.
\end{proof}

\begin{table}[htbp]
\centering
\caption{The first Stirling coefficients $c_j$ in
$\GammaStar(n)\sim\sum c_jn^{-j}$.}
\label{tab:stirling-coefficients}
\small
\begin{tabular}{@{}r l r l@{}}
\toprule
$j$ & $c_j$ & $j$ & $c_j$\\
\midrule
0 & $1$ & 7 & $-534703531/902961561600$\\
1 & $1/12$ & 8 & $-4483131259/86684309913600$\\
2 & $1/288$ & 9 & $432261921612371/514904800886784000$\\
3 & $-139/51840$ & 10 & $6232523202521089/86504006548979712000$\\
4 & $-571/2488320$ & 11 & $-25834629665134204969/13494625021640835072000$\\
5 & $163879/209018880$ & 12 & $-1579029138854919086429/9716130015581401251840000$\\
6 & $5246819/75246796800$ & & \\
\bottomrule
\end{tabular}
\end{table}

\subsection{Divergence versus exact reconstruction}
The Bernoulli asymptotic
\[
 B_{2m}=(-1)^{m-1}\frac{2(2m)!}{(2\pi)^{2m}}\zeta(2m)
\]
shows that the logarithmic coefficients grow factorially.  Hence neither
\cref{eq:log-gammastar-asymp} nor \cref{eq:gammastar-series} is a convergent
power series in $1/n$.  This does not impair the exact result:
\begin{enumerate}
 \item use $n!$ in \cref{eq:exact-real-transseries} and no Stirling expansion
       is present;
 \item use $\GammaStar(n)$ in \cref{eq:exact-factorization} and the
       fluctuation is exact;
 \item use finitely many $c_j$ for a Poincar\'e approximation;
 \item use the Binet integral \cref{eq:binet-laplace} to reconstruct the
       positive-axis Borel sum exactly.
\end{enumerate}
This hierarchy is important when the phrase ``full transseries'' is used:
the complete pole content is convergent, while the standard factorial
fluctuation has the familiar asymptotic status of Stirling's series.

\section{The full mixed transseries and rigorous two-parameter errors}
\label{sec:mixed-transseries}

\subsection{Complex actions and the complete coefficient array}
Define the principal complex actions
\begin{equation}\label{eq:Lambda-def}
 \Lambda_k:=\Log\!\left(\frac{\zeta_k}{\rho}\right)
 =\begin{cases}
 0,&k=0,\\
 \alpha_k+\ii\theta_k,&k>0,\\
 \alpha_{|k|}-\ii\theta_{|k|},&k<0.
 \end{cases}
\end{equation}
Then $u_k=\e^{-\Lambda_k}$ and
\[
 u_k^{n+1}=\e^{-n\Lambda_k}\e^{-\Lambda_k}.
\]

\begin{theorem}[Full elementary-scale transseries]\label{thm:full-transseries}
As $n\to+\infty$ through the positive integers,
\begin{equation}\label{eq:full-transseries}
 \boxed{\displaystyle
 \Fub_n\asympseries
 \mathcal A(n)
 \sum_{k\in\mathbb Z}\e^{-n\Lambda_k}
 \sum_{j\ge0}\frac{C_{j,k}}{n^j},
 \qquad
 C_{j,k}=c_j\e^{-\Lambda_k}
        =c_j\frac{\rho}{\rho+2\pi\ii k}.}
\end{equation}
Equivalently, in conjugate-pair form,
\begin{equation}\label{eq:full-real-transseries}
 \Fub_n\asympseries
 \mathcal A(n)
 \left(\sum_{j\ge0}\frac{c_j}{n^j}\right)
 \left[1+2\sum_{k\ge1}q_k^{n+1}
             \cos((n+1)\theta_k)\right].
\end{equation}
For each fixed pole sector $k$, every algebraic coefficient is obtained from
the principal-sector coefficient by multiplication by $\rho/\zeta_k$.
\end{theorem}

\begin{proof}
Insert the Poincar\'e expansion
$\GammaStar(n)\sim\sum_{j\ge0}c_jn^{-j}$ into the exact factorization
\cref{eq:exact-factorization}.  Next use
$T_n=\sum_k\e^{-n\Lambda_k}\e^{-\Lambda_k}$.  Since the pole sum is
absolutely convergent for each $n\ge1$, the coefficient attached to
$\e^{-n\Lambda_k}n^{-j}$ is exactly
$c_j\e^{-\Lambda_k}$.  Pairing conjugates gives
\cref{eq:full-real-transseries}.
\end{proof}

\begin{corollary}[General coefficient formulae in every sector]
\label{cor:Cjk}
For $j\ge0$ and $k\in\mathbb Z$,
\begin{align}
 C_{j,k}
 &=\frac{\rho}{\zeta_k}\frac1{j!}
   \BellY_j(1!b_1,\ldots,j!b_j),\label{eq:C-bell}\\
 &=\frac{\rho}{\zeta_k}
   \sum_{\substack{m_1+2m_2+\cdots+jm_j=j}}
   \prod_{r=1}^{j}\frac{b_r^{m_r}}{m_r!},
 \label{eq:C-partition}\\
 C_{0,k}&=\frac{\rho}{\zeta_k},\qquad
 C_{j,k}=\frac1j\sum_{r=1}^{j}
       \frac{B_{r+1}}{r+1}C_{j-r,k}.
 \label{eq:C-recurrence}
\end{align}
\end{corollary}

\begin{proof}
Multiply the corresponding formulae of \cref{thm:cj-formulas} by
$\rho/\zeta_k$.  The sector factor is independent of $j$, so the same
triangular recurrence is inherited unchanged.
\end{proof}

\subsection{A rigorous mixed truncation theorem}
Let
\begin{equation}\label{eq:SJ}
 S_J(n):=\sum_{j=0}^{J-1}\frac{c_j}{n^j},
 \qquad J\ge1,
\end{equation}
and retain the pole truncation $T_{n,K}$ from \cref{eq:TK}.  Define
\begin{equation}\label{eq:GJ}
 G_J(n):=
 \begin{cases}
 \displaystyle\frac{1}{\pi^2n},&J=1,\\[2mm]
 \displaystyle\frac{(1+\zeta(J))\Gamma(J)}
 {(2\pi)^{J+1}n^J},&J\ge2.
 \end{cases}
\end{equation}
Here $\zeta(J)$ is the Riemann zeta function.

\begin{theorem}[Two-parameter certified remainder]\label{thm:mixed-error}
For integers $n\ge1$, $J\ge1$, and $K\ge0$,
\begin{equation}\label{eq:mixed-error}
 \left|\frac{\Fub_n}{\mathcal A(n)}-S_J(n)T_{n,K}\right|
 \le G_J(n)(1+B_{n,0})+|S_J(n)|B_{n,K},
\end{equation}
where $B_{n,K}$ may be taken from
\cref{eq:B-exact,eq:B-hurwitz,eq:B-integral}.  Consequently, for fixed
$J,K$,
\begin{equation}\label{eq:mixed-bigO}
 \frac{\Fub_n}{\mathcal A(n)}
 =S_J(n)T_{n,K}+O(n^{-J})
  +O\!\left(q_{K+1}^{n+1}\right).
\end{equation}
\end{theorem}

\begin{proof}
For positive real $n$, the standard scaled-gamma remainder bound
\cite[\S5.11(ii)]{DLMF} gives
\begin{equation}\label{eq:gammastar-bound}
 |\GammaStar(n)-S_J(n)|\le G_J(n).
\end{equation}
For $J\ge2$ this is the positive-axis specialization of the classical bound
\[
 |R_J(z)|\le
 \frac{(1+\zeta(J))\Gamma(J)}
 {2(2\pi)^{J+1}|z|^J}
 \left(1+\min(\sec(\arg z),2\sqrt J)\right);
\]
at $z=n>0$ the last parenthesis is $2$.  For $J=1$, the prescribed
replacement of $1+\zeta(J)$ by $4$ gives $1/(\pi^2n)$.

Using the exact identity $\Fub_n/\mathcal A(n)=\GammaStar(n)T_n$,
\[
 \GammaStar(n)T_n-S_J(n)T_{n,K}
 =\bigl(\GammaStar(n)-S_J(n)\bigr)T_n
  +S_J(n)\bigl(T_n-T_{n,K}\bigr).
\]
Now $|T_n|\le1+B_{n,0}$ and
$|T_n-T_{n,K}|\le B_{n,K}$.  The triangle inequality proves
\cref{eq:mixed-error}; the fixed-truncation asymptotic statement follows.
\end{proof}

\begin{interpretbox}
The two truncation directions behave differently.
\begin{itemize}
 \item Increasing $K$ adds terms of a convergent series and monotonically
       improves the rigorous envelope $B_{n,K}$.
 \item Increasing $J$ traverses a divergent Poincar\'e series.  It improves the
       approximation only until the terms approach their least magnitude;
       the exact Binet integral remains available beyond that point.
\end{itemize}
Thus pole truncation is ordinary convergence, whereas algebraic truncation is
asymptotic optimization.
\end{interpretbox}

\subsection{Numerical behavior}
The next table records the actual absolute relative error
\[
 \left|\frac{n!}{2\rho^{n+1}}\frac{T_{n,K}}{\Fub_n}-1\right|
 =\left|\frac{T_{n,K}}{T_n}-1\right|
\]
for several $n$ and numbers $K$ of retained conjugate pole pairs.

\begin{table}[htbp]
\centering
\caption{Actual relative error of exact-factorial pole truncations.}
\label{tab:pole-errors}
\small
\begin{tabular}{@{}rllll@{}}
\toprule
$n$ & $K=0$ & $K=1$ & $K=2$ & $K=3$\\
\midrule
1  & $4.06845\times 10^{-2}$ & $1.62606\times 10^{-2}$ & $9.98550\times 10^{-3}$ & $7.18241\times 10^{-3}$\\
2  & $9.26902\times 10^{-4}$ & $7.26028\times 10^{-5}$ & $1.75705\times 10^{-5}$ & $6.63865\times 10^{-6}$\\
3  & $2.85346\times 10^{-4}$ & $2.37599\times 10^{-5}$ & $5.80868\times 10^{-6}$ & $2.20191\times 10^{-6}$\\
5  & $2.80687\times 10^{-6}$ & $5.88313\times 10^{-8}$ & $6.03717\times 10^{-9}$ & $1.23133\times 10^{-9}$\\
10 & $5.15518\times 10^{-11}$ & $1.62477\times 10^{-14}$ & $1.34341\times 10^{-16}$ & $4.50497\times 10^{-18}$\\
20 & $1.02594\times 10^{-20}$ & $6.65135\times 10^{-27}$ & $1.03570\times 10^{-30}$ & $1.95587\times 10^{-33}$\\
50 & $1.38135\times 10^{-49}$ & $4.00015\times 10^{-65}$ & $1.28072\times 10^{-73}$ & $5.70913\times 10^{-80}$\\
100& $2.19123\times 10^{-97}$ & $9.02847\times 10^{-128}$& $1.32425\times 10^{-145}$& $2.12149\times 10^{-158}$\\
\bottomrule
\end{tabular}
\end{table}

The nonmonotone-looking comparison between $n=2$ and $n=3$ at $K=0$ is a
phase effect, not a failure of the envelope.  The bounds in
\cref{thm:pole-tail} decrease monotonically with $n$ and $K$ even when the
actual oscillatory error experiences cancellation.

For the mixed elementary approximation, representative relative errors are
shown in \cref{tab:mixed-errors}.  Here $(J,K)$ means $J$ Stirling terms and
$K$ conjugate pole pairs.

\begin{table}[htbp]
\centering
\caption{Actual relative error of selected mixed truncations.}
\label{tab:mixed-errors}
\begingroup
\scriptsize
\setlength{\tabcolsep}{3pt}
\begin{tabular}{@{}rllllll@{}}
\toprule
$n$ & $(1,0)$ & $(3,0)$ & $(6,0)$ & $(3,1)$ & $(6,1)$ & $(10,1)$\\
\midrule
5  & $1.65042\times 10^{-2}$ & $2.40206\times 10^{-5}$ & $2.80967\times 10^{-6}$ & $2.12725\times 10^{-5}$ & $6.16319\times 10^{-8}$ & $5.88588\times 10^{-8}$\\
10 & $8.29596\times 10^{-3}$ & $2.67411\times 10^{-6}$ & $4.08192\times 10^{-11}$& $2.67405\times 10^{-6}$ & $1.07163\times 10^{-11}$& $2.76368\times 10^{-14}$\\
20 & $4.15765\times 10^{-3}$ & $3.34956\times 10^{-7}$ & $6.23885\times 10^{-13}$& $3.34956\times 10^{-7}$ & $6.23885\times 10^{-13}$& $2.26443\times 10^{-18}$\\
50 & $1.66526\times 10^{-3}$ & $2.14490\times 10^{-8}$ & $3.69757\times 10^{-15}$& $2.14490\times 10^{-8}$ & $3.69757\times 10^{-15}$& $3.44988\times 10^{-22}$\\
\bottomrule
\end{tabular}
\endgroup
\end{table}


\section{Nonlinear transseries and multi-index coefficient calculus}
\label{sec:nonlinear}

\subsection{A uniformly small exponential correction}
Write
\begin{equation}\label{eq:Delta-def}
 T_n=1+\Delta_n,
 \qquad
 \Delta_n:=\sum_{k\in\mathbb Z\setminus\{0\}}u_k^{n+1}.
\end{equation}
The proof of \cref{eq:global-relative-bound} gives the stronger absolute-majorant
statement
\begin{equation}\label{eq:Delta-majorant}
 \sum_{k\ne0}|u_k|^{n+1}
 \le\frac{\rho^2}{12}<0.041,
 \qquad n\ge1.
\end{equation}
Therefore Taylor expansions of analytic functions around $T_n=1$ are not
merely formal: they converge absolutely and uniformly in $n\ge1$ whenever
their disk of analyticity has radius greater than $\rho^2/12$.

Let $I:=\mathbb Z\setminus\{0\}$.  A \emph{finite multi-index} is a family
$\bm m=(m_k)_{k\in I}$ of nonnegative integers with finite support.  Put
\[
 |\bm m|:=\sum_{k\in I}m_k,
 \qquad
 \bm m!:=\prod_{k\in I}m_k!,
 \qquad
 u^{\bm m}:=\prod_{k\in I}u_k^{m_k}.
\]
The zero multi-index has $|\bm0|=0$, $\bm0!=1$, and $u^{\bm0}=1$.

\begin{theorem}[Master analytic-transform formula]\label{thm:analytic-transform}
Let $\Phi$ be analytic in $|z-1|<R$ for some
$R>\rho^2/12$.  Then for every $n\ge1$,
\begin{equation}\label{eq:analytic-transform}
 \boxed{\displaystyle
 \Phi(T_n)=
 \sum_{\bm m}
 \frac{\Phi^{(|\bm m|)}(1)}{\bm m!}
 \bigl(u^{\bm m}\bigr)^{n+1}.}
\end{equation}
The sum extends over all finite multi-indices and is absolutely convergent.
It is the complete nonlinear exponential transseries of $\Phi(T_n)$.
\end{theorem}

\begin{proof}
Taylor's theorem gives
\[
 \Phi(1+\Delta_n)=\sum_{r\ge0}\frac{\Phi^{(r)}(1)}{r!}\Delta_n^r.
\]
By \cref{eq:Delta-majorant}, the series converges absolutely.  The multinomial
expansion in countably many absolutely summable variables is
\[
 \Delta_n^r
 =\sum_{\substack{\bm m\\|\bm m|=r}}
   \frac{r!}{\bm m!}
   \prod_{k\in I}\bigl(u_k^{n+1}\bigr)^{m_k}
 =\sum_{|\bm m|=r}\frac{r!}{\bm m!}
   \bigl(u^{\bm m}\bigr)^{n+1}.
\]
Substitution cancels $r!$ and proves \cref{eq:analytic-transform}.  Absolute
convergence follows by repeating the same calculation with absolute values
and using a slightly smaller circle inside the disk of analyticity.
\end{proof}

\subsection{Powers, reciprocals, and logarithms}
Taking $\Phi(z)=z^\sigma$, with the branch analytic in $|z-1|<1$, gives
\begin{equation}\label{eq:T-power}
 T_n^\sigma
 =\sum_{\bm m}
   \frac{\falling{\sigma}{|\bm m|}}{\bm m!}
   \bigl(u^{\bm m}\bigr)^{n+1},
 \qquad \sigma\in\mathbb C,
\end{equation}
where
\(
 \falling{\sigma}{r}=\sigma(\sigma-1)\cdots(\sigma-r+1)
\).
Two important special cases are
\begin{align}
 \frac1{T_n}
 &=\sum_{\bm m}(-1)^{|\bm m|}
       \frac{|\bm m|!}{\bm m!}
       \bigl(u^{\bm m}\bigr)^{n+1},
 \label{eq:T-reciprocal}\\
 \log T_n
 &=\sum_{\bm m\ne\bm0}(-1)^{|\bm m|-1}
       \frac{(|\bm m|-1)!}{\bm m!}
       \bigl(u^{\bm m}\bigr)^{n+1}.
 \label{eq:T-log}
\end{align}
The coefficient in \cref{eq:T-log} is the connected, or cumulant, version of
the multinomial coefficient.  Products of distinct pole sectors and repeated
powers of a single sector appear automatically.

\subsection{Arbitrary powers of the Fubini numbers}
For $\sigma\in\mathbb C$, define the generalized Stirling coefficients
$d_j(\sigma)$ by
\begin{equation}\label{eq:d-sigma}
 \exp\!\left(\sigma\sum_{r\ge1}b_rt^r\right)
 =\sum_{j\ge0}d_j(\sigma)t^j.
\end{equation}
The same Bell-polynomial argument as before gives
\begin{align}
 d_j(\sigma)
 &=\frac1{j!}\BellY_j\bigl(\sigma1!b_1,\ldots,
                             \sigma j!b_j\bigr),
 \label{eq:d-bell}\\
 d_0(\sigma)&=1,
 \qquad
 d_j(\sigma)=\frac{\sigma}{j}\sum_{r=1}^{j}
     \frac{B_{r+1}}{r+1}d_{j-r}(\sigma).
 \label{eq:d-recurrence}
\end{align}
Thus
\(
 \GammaStar(n)^\sigma\sim\sum_{j\ge0}d_j(\sigma)n^{-j}
\).
Combining this with \cref{eq:T-power} yields the full coefficient calculus
for arbitrary powers.

\begin{corollary}[Full transseries for $\Fub_n^\sigma$]
\label{cor:F-power}
For any fixed $\sigma\in\mathbb C$,
\begin{equation}\label{eq:F-power}
 \Fub_n^\sigma\asympseries
 \mathcal A(n)^\sigma
 \sum_{\bm m}
 \frac{\falling{\sigma}{|\bm m|}}{\bm m!}
 \bigl(u^{\bm m}\bigr)^{n+1}
 \sum_{j\ge0}\frac{d_j(\sigma)}{n^j}.
\end{equation}
For positive real $\Fub_n$ the principal real power is understood; the
right-hand side supplies its analytic continuation in $\sigma$.
\end{corollary}

\subsection{The complete logarithmic transseries}
Taking the logarithm of \cref{eq:exact-factorization} gives an exact additive
decomposition:
\begin{align}
 \log\Fub_n
 ={}&(n+\tfrac12)\log n-n-(n+1)\log\rho
      +\tfrac12\log(2\pi)-\log2\nonumber\\
 &+\log\GammaStar(n)+\log T_n.
 \label{eq:log-exact-decomp}
\end{align}
Inserting the Binet integral and the convergent multi-index sum yields an
entirely exact representation.

\begin{theorem}[Exact and asymptotic logarithmic forms]\label{thm:log-F}
For every $n\ge1$,
\begin{align}
 \log\Fub_n
 ={}&(n+\tfrac12)\log n-n-(n+1)\log\rho
      +\tfrac12\log(2\pi)-\log2\nonumber\\
 &+\int_0^\infty \e^{-nt}h(t)\dd t
 +\sum_{\bm m\ne\bm0}(-1)^{|\bm m|-1}
       \frac{(|\bm m|-1)!}{\bm m!}
       \bigl(u^{\bm m}\bigr)^{n+1}.
 \label{eq:log-F-exact}
\end{align}
The integral and multi-index sum converge.  Its all-orders Poincar\'e form is
\begin{align}
 \log\Fub_n\asympseries{}&(n+\tfrac12)\log n-n-(n+1)\log\rho
      +\tfrac12\log(2\pi)-\log2\nonumber\\
 &+\sum_{r\ge1}\frac{B_{r+1}}{r(r+1)n^r}
 +\sum_{\bm m\ne\bm0}(-1)^{|\bm m|-1}
       \frac{(|\bm m|-1)!}{\bm m!}
       \bigl(u^{\bm m}\bigr)^{n+1}.
 \label{eq:log-F-asymp}
\end{align}
\end{theorem}

\begin{proof}
Equation \cref{eq:log-exact-decomp}, Binet's formula
\cref{eq:binet-laplace}, and the absolutely convergent identity
\cref{eq:T-log} prove \cref{eq:log-F-exact}.  Replacing the integral by its
Watson expansion \cref{eq:log-gammastar-asymp} gives
\cref{eq:log-F-asymp}.
\end{proof}

The first nonlinear exponential terms are
\begin{align*}
 \log T_n={}&
 \sum_{k\ne0}u_k^{n+1}
 -\frac12\sum_{k,\ell\ne0}(u_ku_\ell)^{n+1}
 +\frac13\sum_{k,\ell,m\ne0}(u_ku_\ell u_m)^{n+1}-\cdots.
\end{align*}
The multi-index form is preferable because it collects repeated products and
gives their exact multinomial coefficients.

\subsection{The ratio of consecutive Fubini numbers}
The exact factorial factors cancel almost completely:
\begin{equation}\label{eq:ratio-exact}
 \mathcal R_n:=\frac{n\Fub_{n-1}}{\Fub_n}
 =\rho\frac{T_{n-1}}{T_n},
 \qquad n\ge1.
\end{equation}
For an absolutely convergent multi-index expansion of the numerator we take
$n\ge2$.  In terms of $y_k=u_k^{n+1}$,
\[
 T_n=1+\sum_{k\ne0}y_k,
 \qquad
 T_{n-1}=1+\sum_{k\ne0}u_k^{-1}y_k.
\]

\begin{theorem}[General ratio coefficients]\label{thm:ratio}
For $n\ge2$,
\begin{equation}\label{eq:ratio-multi}
 \frac{\mathcal R_n}{\rho}
 =1+\sum_{\bm m\ne\bm0}r_{\bm m}
      \bigl(u^{\bm m}\bigr)^{n+1},
\end{equation}
where, with $M=|\bm m|$,
\begin{equation}\label{eq:ratio-coeff}
 \boxed{\displaystyle
 r_{\bm m}=(-1)^{M-1}\frac{(M-1)!}{\bm m!}
 \left(\sum_{k\ne0}m_ku_k^{-1}-M\right).}
\end{equation}
The series is absolutely convergent.  In particular,
\begin{align}
 \frac{\mathcal R_n}{\rho}
 ={}&1+\sum_{k\ne0}(u_k^{-1}-1)u_k^{n+1}
 +O\!\left(q_2^n+q_1^{2n}\right)\nonumber\\
 ={}&1+2q_1^n\left[
 \cos(n\theta_1)-q_1\cos((n+1)\theta_1)
 \right]
 +O\!\left(q_2^n+q_1^{2n}\right).
 \label{eq:ratio-first}
\end{align}
Consequently
\(
 \Fub_n/\Fub_{n-1}\sim n/\rho
\), with exponentially small rather than algebraic relative corrections.
\end{theorem}

\begin{proof}
Multiply
\[
 \left(1+\sum_{k\ne0}u_k^{-1}y_k\right)
 \left(1+\sum_{k\ne0}y_k\right)^{-1}.
\]
For a nonzero multi-index $\bm m$ of total degree $M$, the contribution from
the constant numerator is
$(-1)^MM!/\bm m!$.  Choosing the term $u_k^{-1}y_k$ from the numerator leaves
multi-index $\bm m-\bm e_k$ in the reciprocal and contributes
\[
 u_k^{-1}(-1)^{M-1}\frac{(M-1)!}{(\bm m-\bm e_k)!}
 =(-1)^{M-1}\frac{(M-1)!}{\bm m!}m_ku_k^{-1}.
\]
Summing over $k$ and combining the constant-numerator term gives
\cref{eq:ratio-coeff}.  Absolute convergence follows from
$\sum_{k\ne0}|u_k|^n<\infty$ for $n\ge2$ and from
$\sum|u_k|^{n+1}<1$.  Keeping only the linear multi-indices and pairing
$k=\pm1$ gives \cref{eq:ratio-first}; the next envelopes come from the
$k=\pm2$ linear terms and quadratic products of the first pair.
\end{proof}

\section{Weighted Fubini polynomials and geometric moments}
\label{sec:weighted}

\subsection{Definition and generating functions}
For $x>0$, define the Fubini polynomial
\begin{equation}\label{eq:Fubini-poly-def}
 \Fub_n(x):=\sum_{j=0}^{n}j!\StirlingII{n}{j}x^j.
\end{equation}
The same labelled-block argument as in \cref{thm:egf} gives
\begin{equation}\label{eq:Fubini-poly-egf}
 \sum_{n\ge0}\Fub_n(x)\frac{z^n}{n!}
 =\frac{1}{1-x(\e^z-1)}
 =\frac{1}{1+x-x\e^z}.
\end{equation}
Put
\begin{equation}\label{eq:q-rhox}
 q:=\frac{x}{1+x},
 \qquad
 \rho_x:=-\log q=\log\!\left(1+\frac1x\right).
\end{equation}
Then
\begin{equation}\label{eq:weighted-moment}
 \Fub_n(x)=(1-q)\sum_{m\ge0}q^m m^n.
\end{equation}
Thus $\Fub_n(x)$ is the $n$th moment of the geometric distribution
$\Pr(G_x=m)=(1-q)q^m$.  For $n\ge1$,
\begin{equation}\label{eq:weighted-polylog}
 \Fub_n(x)=(1-q)\Li_{-n}(q).
\end{equation}
The Eulerian form is
\begin{equation}\label{eq:weighted-eulerian}
 \Fub_n(x)=x(1+x)^{n-1}
 A_n\!\left(\frac{x}{1+x}\right),
 \qquad n\ge1.
\end{equation}

\subsection{Exact poles and the weighted transseries}
The poles of \cref{eq:Fubini-poly-egf} are
\[
 \zeta_k(x)=\rho_x+2\pi\ii k,
 \qquad k\in\mathbb Z,
\]
and the residue at every pole is $-1/(1+x)$.  Repeating the coth decomposition
or Poisson summation gives the following master extension.

\begin{theorem}[Full weighted pole formula]\label{thm:weighted-poles}
For $x>0$ and $n\ge1$,
\begin{equation}\label{eq:weighted-poles}
 \boxed{\displaystyle
 \Fub_n(x)=\frac{n!}{1+x}
 \sum_{k\in\mathbb Z}
 \frac{1}{(\rho_x+2\pi\ii k)^{n+1}}.}
\end{equation}
Let
\[
 u_k(x):=\frac{\rho_x}{\rho_x+2\pi\ii k},
 \qquad
 T_n(x):=\sum_{k\in\mathbb Z}u_k(x)^{n+1}.
\]
Then
\begin{equation}\label{eq:weighted-exact-factor}
 \Fub_n(x)=\mathcal A_x(n)\GammaStar(n)T_n(x),
 \qquad
 \mathcal A_x(n):=
 \frac{\sqrt{2\pi n}}{(1+x)\rho_x}
 \left(\frac{n}{\e\rho_x}\right)^n.
\end{equation}
Its full elementary-scale coefficients are
\begin{equation}\label{eq:weighted-Cjk}
 C_{j,k}(x)=c_j\frac{\rho_x}{\rho_x+2\pi\ii k}.
\end{equation}
In particular,
\begin{equation}\label{eq:weighted-leading}
 \Fub_n(x)\sim\frac{n!}{(1+x)\rho_x^{n+1}}
 \qquad(n\to\infty,\text{ fixed }x>0).
\end{equation}
\end{theorem}

\begin{proof}
Factor the EGF as
\[
 \frac{1}{1+x-x\e^z}
 =\frac{1}{1+x}\frac{1}{1-q\e^z}.
\]
With $w=z-\rho_x$, one has $q\e^z=\e^w$, hence
\[
 \frac{1}{1+x-x\e^z}
 =\frac{1}{2(1+x)}-
  \frac{1}{2(1+x)}\coth\!\left(\frac w2\right).
\]
The coefficient extraction of \cref{thm:mittag,thm:exact-poles} now applies
with residue $-1/(1+x)$ and $\rho$ replaced by $\rho_x$, proving
\cref{eq:weighted-poles}.  The remaining statements follow from
\cref{eq:factorial-scaled} exactly as before.
\end{proof}

\begin{corollary}[Uniformity on compact weight intervals]
Fix $0<a\le x\le b<\infty$.  For every fixed $K$ and $J$, the weighted
analogues of \cref{eq:mixed-bigO} hold uniformly in $x\in[a,b]$.
\end{corollary}

\begin{proof}
On $[a,b]$, $\rho_x$ stays in a compact subinterval of $(0,\infty)$.  Hence
$|u_k(x)|$ is uniformly bounded away from one for every fixed $k\ne0$, and
the tails are uniformly dominated by a constant multiple of
$\sum_{k>K}k^{-n-1}$.  The scaled-gamma remainder is independent of $x$.
\end{proof}

\begin{remark}[Nonuniform endpoint regimes]
The fixed-$x$ transseries is not uniform as $x\downarrow0$.  In that limit
$\rho_x\to\infty$ and
$|u_1(x)|\to1$, so many pole sectors become comparable before the eventual
$n\to\infty$ hierarchy separates them.  By contrast, as $x\to\infty$,
$\rho_x\sim x^{-1}$ and the nonreal sectors are extremely suppressed.  A
double-scaling theory in which $n$ and $x$ vary together is therefore a
separate asymptotic problem rather than a formal substitution into the
fixed-$x$ result.
\end{remark}

\subsection{Weighted cumulants and Bell-polynomial reconstruction}
The cumulant generating function of the geometric law in
\cref{eq:weighted-moment} is
\[
 K_x(z)=\log(1-q)-\log(1-q\e^z).
\]
Thus
\begin{equation}\label{eq:weighted-cumulants}
 \kappa_r(x)=\Li_{1-r}(q),
 \qquad q=\frac{x}{1+x},
\end{equation}
and
\begin{equation}\label{eq:weighted-bell}
 \Fub_n(x)=\BellY_n(\kappa_1(x),\ldots,\kappa_n(x)).
\end{equation}
This identity supplies an exact coefficient bridge between ordered set
partitions, geometric moments, and the Bell-polynomial calculus used for the
Stirling fluctuation.  The two appearances of Bell polynomials play different
roles: \cref{eq:weighted-bell} converts probabilistic cumulants into exact
moments, whereas \cref{eq:c-bell} exponentiates the asymptotic logarithm of a
gamma factor.


\section{A general meromorphic-EGF principle and higher-order poles}
\label{sec:meromorphic}

\subsection{Coefficient contributions of arbitrary principal parts}
The Fubini case is a particularly transparent instance of a general theorem.

\begin{theorem}[Meromorphic exponential generating functions]
\label{thm:meromorphic-egf}
Let
\[
 A(z)=\sum_{n\ge0}a_n\frac{z^n}{n!}
\]
be meromorphic in a neighborhood of the closed disk $|z|\le R$, with no pole
on $|z|=R$.  Let its poles in the disk be $\xi_1,\ldots,\xi_M$, and write the
principal part at $\xi_\ell$ as
\[
 \sum_{s=1}^{d_\ell}a_{\ell,-s}(z-\xi_\ell)^{-s}.
\]
Then
\begin{equation}\label{eq:meromorphic-coeff}
 a_n=n!\sum_{\ell=1}^{M}\sum_{s=1}^{d_\ell}
 (-1)^s a_{\ell,-s}
 \binom{n+s-1}{s-1}\xi_\ell^{-n-s}+E_n,
\end{equation}
where, after subtracting all principal parts, the analytic remainder $H$ obeys
\begin{equation}\label{eq:meromorphic-error}
 |E_n|\le n!R^{-n}\max_{|z|=R}|H(z)|.
\end{equation}
In particular, a simple pole of residue $r_\ell$ contributes exactly
\[
 -n!r_\ell\xi_\ell^{-n-1},
\]
whereas a pole of order $d$ contributes an exponential factor
$\xi^{-n}$ multiplied by a polynomial in $n$ of degree at most $d-1$.
\end{theorem}

\begin{proof}
Subtract the principal parts and call the result $H(z)$.  It is analytic on a
neighborhood of the closed disk.  For a single principal-part term,
\[
 (z-\xi)^{-s}
 =(-\xi)^{-s}(1-z/\xi)^{-s}
 =(-1)^s\sum_{n\ge0}
   \binom{n+s-1}{s-1}\xi^{-n-s}z^n.
\]
Summing these coefficient contributions and multiplying by $n!$ gives the
main term in \cref{eq:meromorphic-coeff}.  Cauchy's coefficient estimate on
$|z|=R$ gives \cref{eq:meromorphic-error}.  The binomial coefficient is a
polynomial in $n$ of degree $s-1$, proving the final statement.
\end{proof}

\begin{remark}
For a function with infinitely many poles one applies
\cref{thm:meromorphic-egf} on expanding disks.  If a global Mittag--Leffler
decomposition is available, as in \cref{thm:mittag}, the limiting sum may be
exact.  The theorem also explains the sector ordering: poles are ranked by
modulus, while equal-modulus conjugates must be retained together for a real
approximation.
\end{remark}

\subsection{Higher-order Fubini numbers}
For an integer $r\ge1$, define the order-$r$ Fubini numbers by
\begin{equation}\label{eq:higher-fubini-def}
 \frac{1}{(2-\e^z)^r}
 =\sum_{n\ge0}\Fub_n^{\langle r\rangle}\frac{z^n}{n!}.
\end{equation}
Every pole now has order $r$.  Introduce the generalized Bernoulli, or
N\o rlund, numbers by
\begin{equation}\label{eq:norlund-def}
 \left(\frac{t}{\e^t-1}\right)^r
 =\sum_{j\ge0}B_j^{(r)}\frac{t^j}{j!}.
\end{equation}
For example,
\[
 B_0^{(r)}=1,
 \qquad B_1^{(r)}=-\frac r2,
 \qquad B_2^{(r)}=\frac{r(3r-1)}{12}.
\]

\begin{theorem}[Exact pole transseries for higher order]
\label{thm:higher-fubini}
For integers $r\ge1$ and $n\ge1$,
\begin{equation}\label{eq:higher-fubini-poles}
 \boxed{\displaystyle
 \Fub_n^{\langle r\rangle}
 =\frac{n!}{2^r}\sum_{k\in\mathbb Z}
   \sum_{j=0}^{r-1}
   \frac{(-1)^jB_j^{(r)}}{j!}
   \binom{n+r-j-1}{r-j-1}
   \zeta_k^{-n-r+j}.}
\end{equation}
The sums are absolutely convergent.  In particular,
\begin{equation}\label{eq:higher-leading}
 \Fub_n^{\langle r\rangle}
 \sim\frac{n!}{2^r}
 \binom{n+r-1}{r-1}\rho^{-n-r}.
\end{equation}
Every pole sector has a polynomial amplitude of degree $r-1$, whose complete
coefficients are the generalized Bernoulli numbers in
\cref{eq:higher-fubini-poles}.
\end{theorem}

\begin{proof}
Near $z=\zeta_k$, put $w=z-\zeta_k$.  Since $\e^{\zeta_k}=2$,
\begin{align*}
 \frac{1}{(2-\e^z)^r}
 &=\frac{1}{2^r(1-\e^w)^r}\\
 &=\frac{(-1)^r}{2^r}w^{-r}
   \left(\frac{w}{\e^w-1}\right)^r\\
 &=\frac{(-1)^r}{2^r}
   \sum_{j\ge0}\frac{B_j^{(r)}}{j!}w^{j-r}.
\end{align*}
The principal part consists of $0\le j\le r-1$.  Apply
\cref{thm:meromorphic-egf} to each principal-part term.  If
$s=r-j$, then the sign from coefficient extraction is $(-1)^s$; combined
with $(-1)^r$ it is $(-1)^j$.  This gives exactly the summand in
\cref{eq:higher-fubini-poles}.

For completeness, exactness rather than merely an asymptotic sum can be seen
without any global convergence subtlety.  For $a>1$ and $n\ge1$,
\[
 n!\Coeff{z}{n}{\frac{1}{a-\e^z}}
 =\frac{n!}{a}\sum_{k\in\mathbb Z}
    (\log a+2\pi\ii k)^{-n-1}.
\]
Moreover,
\[
 (a-\e^z)^{-r}
 =\frac{(-1)^{r-1}}{(r-1)!}
   \frac{\partial^{r-1}}{\partial a^{r-1}}(a-\e^z)^{-1}.
\]
The bilateral series and its first $r-1$ parameter derivatives converge
locally uniformly for $n\ge1$, so differentiation followed by $a=2$ proves
that the sum of all pole contributions is exact.  Finally the unique dominant
term is $k=0,j=0$, which gives \cref{eq:higher-leading}.
\end{proof}

\begin{example}[Order two]
Since $B_0^{(2)}=1$ and $B_1^{(2)}=-1$,
\begin{equation}\label{eq:order-two}
 \Fub_n^{\langle2\rangle}
 =\frac{n!}{4}\sum_{k\in\mathbb Z}
 \left[(n+1)\zeta_k^{-n-2}+\zeta_k^{-n-1}\right].
\end{equation}
The first values are
\(
 1,2,8,44,308,2612,25988,\ldots
\).
The linear polynomial $n+1$ is the characteristic signature of a double
pole.
\end{example}

\subsection{What the general theorem says about ``full'' expansions}
The preceding results separate three sources of coefficients that are often
mixed together:
\begin{enumerate}
 \item the \emph{principal-part coefficients} at a generating-function pole,
       which determine the polynomial sector amplitude;
 \item the \emph{pole location}, which determines the exponential action and
       phase;
 \item any \emph{external normalization}, such as replacing $n!$ by an
       elementary Stirling scale, which supplies a common algebraic series.
\end{enumerate}
For ordinary Fubini numbers the first item is only the residue $-1/2$, so the
native pole sectors have no internal $1/n$ series.  Formula
\cref{eq:full-transseries} contains inverse powers only because of the third
item.  Formula \cref{eq:higher-fubini-poles} shows explicitly how the answer
changes when the first item becomes nontrivial.

\section{Derived asymptotic quantities}\label{sec:derived}

\subsection{Root asymptotics}
Dividing \cref{eq:log-F-exact} by $n$ gives a convenient form for the root
growth.  In particular,
\begin{align}
 \Fub_n^{1/n}
 ={}&\frac{n}{\e\rho}
 \exp\!\left[
  \frac{\frac12\log(2\pi n)-\log(2\rho)}{n}
  +\frac{1}{12n^2}-\frac{1}{360n^4}+\frac{1}{1260n^6}
  +O(n^{-8})
 \right]\nonumber\\
 &\times
 \exp\!\left[
  \frac{2q_1^{n+1}\cos((n+1)\theta_1)+O(q_2^{n+1}+q_1^{2n+2})}{n}
 \right].
 \label{eq:nth-root}
\end{align}
The first line is the algebraic expansion; the second is the exponentially
small oscillatory correction.  At leading order,
\[
 \Fub_n^{1/n}\sim\frac{n}{\e\log2}.
\]

\subsection{Successive quotients}
The ratio theorem gives
\begin{equation}\label{eq:successive-quotient}
 \frac{\Fub_n}{\Fub_{n-1}}
 =\frac{n}{\rho}
 \left\{1-2q_1^n\left[
 \cos(n\theta_1)-q_1\cos((n+1)\theta_1)
 \right]+O(q_2^n+q_1^{2n})\right\}.
\end{equation}
There are no algebraic $1/n$ corrections.  This cancellation is exact:
$n!/(n-1)!=n$, and the common power of $\rho$ contributes $1/\rho$.
Only the exponentially small pole corrections survive.

\subsection{Logarithmic convexity at large order}
For $n\ge1$, the geometric-moment representation implies log-convexity
directly by Cauchy--Schwarz:
\[
 \Fub_n^2=\mathbb E[G^n]^2
 \le\mathbb E[G^{n-1}]\mathbb E[G^{n+1}]
 =\Fub_{n-1}\Fub_{n+1}.
\]
The transseries refines the asymptotic size of the log-convexity quotient.  Set
\[
 Q_n:=\frac{\Fub_{n-1}\Fub_{n+1}}{\Fub_n^2}.
\]
The factorial and dominant-pole factors give
\(
 Q_n=(n+1)T_{n-1}T_{n+1}/(nT_n^2)
\), hence
\begin{equation}\label{eq:logconvex-ratio}
 Q_n=1+\frac1n+O(q_1^{n}).
\end{equation}
A complete multi-index expansion follows from
\cref{thm:analytic-transform} after adjoining the shifted sector variables
$u_k^{n}$ and $u_k^{n+2}$.

\section{Computation, verification, and implementation}\label{sec:computation}

\subsection{Which formula should be used?}
The available formulae serve different computational regimes.
\begin{enumerate}
 \item For small exact $n$, use the Stirling sum
       \cref{thm:stirling-sum} or the binomial recurrence
       \cref{cor:recurrence}; both use exact integers.
 \item For high-precision asymptotics with moderate or large $n$, evaluate a
       few conjugate pole pairs in polar form and keep $n!$ through
       $\log\Gamma(n+1)$.
 \item For an elementary closed approximation, use $S_J(n)T_{n,K}$ and the
       certified error \cref{eq:mixed-error}.
 \item For extremely high algebraic accuracy, use the Binet integral rather
       than forcing the divergent Stirling series past its least term.
\end{enumerate}

The pole-pair form avoids cancellation between huge complex quantities:
\[
 \log\!\left(\frac{n!}{2\rho^{n+1}}\right)
 =\log\Gamma(n+1)-\log2-(n+1)\log\rho,
\]
followed by multiplication by the real quantity $T_{n,K}$.

\subsection{A priori choice of the number of pole pairs}
Let
\[
 \eta:=1-\frac{\rho^2}{12}=0.959962248840149\ldots.
\]
By \cref{eq:global-relative-bound}, $|T_n|\ge\eta$.  Hence the relative error
of the exact-factorial $K$-pair approximation satisfies
\[
 \left|\frac{T_{n,K}}{T_n}-1\right|\le\frac{B_{n,K}}{\eta}.
\]
For a target relative tolerance $\varepsilon$, it is therefore enough to
choose an integer $K\ge1$ such that
\begin{equation}\label{eq:choose-K}
 \frac{2}{n}\left(\frac{\rho}{2\pi}\right)^{n+1}K^{-n}
 \le\eta\varepsilon.
\end{equation}
Equivalently, one may take
\begin{equation}\label{eq:choose-K-solved}
 K\ge
 \left[
 \frac{2}{n\eta\varepsilon}
 \left(\frac{\rho}{2\pi}\right)^{n+1}
 \right]^{1/n}
\end{equation}
and round the right-hand side upward.  For an absolute tolerance in the
normalized factor $T_n$, simply omit $\eta$.  These bounds are intentionally
simple and conservative; the exact majorant
$B_{n,K}=2\sum_{k>K}q_k^{n+1}$ or the Hurwitz-zeta bound is usually sharper.

\subsection{Wolfram Language reference implementation}
The following definitions compute exact Fubini numbers, pole approximations,
Stirling coefficients, and certified pole-tail envelopes.

\begin{lstlisting}[language=Mathematica,caption={Reference implementation in Wolfram Language.}]
ClearAll[fubiniExact, fubiniPole, stirlingCoefficient, poleTailBound];

fubiniExact[n_Integer?NonNegative] :=
  Sum[k! StirlingS2[n, k], {k, 0, n}];

fubiniPole[n_Integer?Positive, K_Integer?NonNegative,
   wp_Integer : 80] := Module[{rho = N[Log[2], wp]},
  N[n!/2 Sum[(rho + 2 Pi I k)^(-n - 1), {k, -K, K}], wp]
  ];

stirlingCoefficient[0] = 1;
stirlingCoefficient[j_Integer?Positive] :=
  stirlingCoefficient[j] = Together[
    1/j Sum[
      BernoulliB[r + 1]/(r + 1) stirlingCoefficient[j - r],
      {r, 1, j}]
    ];

poleTailBound[n_Integer?Positive, K_Integer?NonNegative] :=
  2 (Log[2]/(2 Pi))^(n + 1) HurwitzZeta[n + 1, K + 1];

Table[fubiniExact[n], {n, 0, 10}]
Table[stirlingCoefficient[j], {j, 0, 12}]
\end{lstlisting}

For a real-valued and phase-stable pole calculation one may replace the
bilateral complex sum with

\begin{lstlisting}[language=Mathematica]
fubiniPoleReal[n_Integer?Positive, K_Integer?NonNegative,
   wp_Integer : 80] := Module[{rho, term},
  rho = N[Log[2], wp];
  term[k_] := Module[{z = rho + 2 Pi I k},
    (rho/Abs[z])^(n + 1) Cos[(n + 1) Arg[z]]
    ];
  N[n!/(2 rho^(n + 1)) (1 + 2 Sum[term[k], {k, 1, K}]), wp]
  ];
\end{lstlisting}

\subsection{Internal consistency checks}
Several identities provide independent tests of an implementation:
\begin{align*}
 \Fub_n
 &=\sum_{j=0}^{n}j!\StirlingII{n}{j},\\
 &=\sum_{j=0}^{n-1}\binom{n}{j}\Fub_j,\\
 &=\frac12\sum_{m\ge0}\frac{m^n}{2^m},\\
 &=\frac{n!}{2}\sum_{k\in\mathbb Z}\zeta_k^{-n-1},\\
 &=2^{n-1}A_n(1/2)\qquad(n\ge1).
\end{align*}
The first two are finite exact integer computations, the third and fourth are
independent rapidly convergent numerical sums, and the fifth uses a different
combinatorial triangle.  Agreement among all five is a strong guard against
normalization, indexing, and $n=0$ errors.

\begin{table}[htbp]
\centering
\caption{Exact initial values used in the numerical checks.}
\label{tab:initial-values}
\begin{tabular}{@{}r r@{\qquad}r r@{}}
\toprule
$n$ & $\Fub_n$ & $n$ & $\Fub_n$\\
\midrule
0 & 1 & 6 & 4683\\
1 & 1 & 7 & 47293\\
2 & 3 & 8 & 545835\\
3 & 13 & 9 & 7087261\\
4 & 75 & 10 & 102247563\\
5 & 541 & 11 & 1622632573\\
\bottomrule
\end{tabular}
\end{table}


\section{Literature context and relation to the linked research drafts}
\label{sec:literature}
The numbers studied here have circulated under several names.  Gross's
``preferential arrangements'' terminology and Good's study of rankings with
ties are standard early modern references
\cite{Gross1962,Good1975}; Comtet popularized the name Fubini numbers in the
coefficient-calculus literature \cite{Comtet1974}.  Good's Theorem~4
already records the exact reciprocal-pole/cosine expansion underlying
\cref{thm:exact-poles,thm:exact-transseries}; the present organization recasts
that identity as a convergent pole-lattice transseries and develops its
all-orders coefficient, remainder, and nonlinear-transform calculus.
Their combinatorial and analytic identities are also catalogued as OEIS
A000670 \cite{OEIS}.  The appearance of Fubini numbers as coefficients in
asymptotic expansions of
hypercube resistance was developed by Pippenger \cite{Pippenger2010}; barred
preferential arrangements give one of many related higher-order families
\cite{Ahlbach2013}.

The leading-pole estimate is a direct instance of meromorphic singularity
analysis, for which \cite{FlajoletSedgewick2009} provides a broad systematic
framework.  What is emphasized in the present article is the \emph{complete}
pole lattice and its exact Mittag--Leffler sum, together with the resulting
transseries algebra.  The Stirling fluctuation and its bounds are standard
gamma-function asymptotics; useful general sources include
\cite{DLMF,Olver1997,WhittakerWatson1927,ParisKaminski2001}.  The distinction
between a divergent Poincar\'e series and its summation is classical
\cite{Hardy1949}.

The notation and organization deliberately align with the two linked
research areas in the ProveIt repository.  The first develops Stirling,
Bell-polynomial, Fa\`a di Bruno, and inversion coefficient calculus
\cite{ProveItCoefficient}; the second organizes formal series and transseries
by their asymptotic scales and algebraic operations \cite{ProveItTransseries}.
The present paper supplies a worked synthesis: the EGF pole lattice provides
the transmonomials, complete Bell polynomials provide all factorial
coefficients, and multi-index multinomial calculus provides nonlinear sector
products.  No claim is made here that the analytic pole and remainder
theorems have been machine-checked in Lean.

\section{Conclusions and further directions}\label{sec:conclusion}
The Fubini numbers admit a stronger description than the customary leading
asymptotic formula suggests.  In exact factorial normalization,
\[
 \frac{2(\log2)^{n+1}}{n!}\Fub_n
 =\sum_{k\in\mathbb Z}
  \left(\frac{\log2}{\log2+2\pi\ii k}\right)^{n+1},
\]
so the complete exponentially small structure is an absolutely convergent
sum indexed by every EGF pole.  The conjugate-pair form makes each real action
and oscillatory phase explicit.  Because the poles are simple, their native
amplitudes are constants.

Replacing $n!$ by its elementary Stirling scale introduces a common
inverse-power series.  Its coefficients are
\[
 c_j=\frac1{j!}\BellY_j(1!b_1,\ldots,j!b_j),
 \qquad
 b_r=\frac{B_{r+1}}{r(r+1)},
\]
and the complete mixed coefficient array is the rank-one factorization
\[
 C_{j,k}=c_j\frac{\log2}{\log2+2\pi\ii k}.
\]
This factorization is the central all-orders result: the $j$-dependence is
universal across sectors, and the $k$-dependence is a single pole multiplier.
The two-parameter theorem \cref{thm:mixed-error} turns any finite rectangle in
this coefficient array into a certified approximation.

The same framework remains exact under analytic nonlinear operations on the
normalized pole sum, yielding closed multi-index coefficients for powers,
reciprocals, logarithms, and ratios.  It extends uniformly to positive weights
on compact intervals and, after replacing simple poles by multiple poles, to
higher-order Fubini families with generalized-Bernoulli polynomial
amplitudes.

Several directions are natural continuations.
\begin{enumerate}
 \item \textbf{Uniform small-weight scaling.}  For $x\downarrow0$,
       $\rho_x\to\infty$ and the fixed-$x$ sector hierarchy loses uniformity.
       Regimes such as $n/\rho_x$ fixed should lead to a different continuum
       approximation to the pole lattice.
 \item \textbf{Complex order and Stokes geometry.}  The continuous-order
       formula \cref{eq:continuous-poles} defines sectors for complex $\nu$.
       Their dominance ordering changes with $\arg\nu$, giving a natural
       Stokes and anti-Stokes diagram even though the positive real sequence
       has a globally convergent pole sum.
 \item \textbf{Hyperasymptotics of the common gamma factor.}  The Borel
       singularities at $\pm2\pi\ii m$ can be resolved with terminant
       functions.  Combining that hyperasymptotic layer with the exact Fubini
       pole lattice would produce a two-geometry resurgent description.
 \item \textbf{Higher-order and barred arrangements.}  Formula
       \cref{eq:higher-fubini-poles} suggests a general catalogue in which
       combinatorial constructions are classified by the orders and residues
       of their meromorphic EGF poles.
 \item \textbf{Formal verification.}  The Stirling-transform EGF, Bell
       coefficient recurrence, finite pole truncations, and algebraic
       multi-index identities are plausible first targets.  The analytic
       Mittag--Leffler, Poisson, and remainder layers require a developed
       complex-analysis interface.
\end{enumerate}

\appendix

\section{Explicit derivation of the first algebraic coefficients}
\label{app:first-coefficients}
The logarithmic coefficients begin
\[
 b_1=\frac1{12},\qquad b_2=0,\qquad
 b_3=-\frac1{360},\qquad b_4=0,\qquad
 b_5=\frac1{1260},\ldots.
\]
Expanding $\exp(\sum b_rt^r)$ by partitions of $j$ gives
\begin{align*}
 c_0&=1,\\
 c_1&=b_1=\frac1{12},\\
 c_2&=b_2+\frac{b_1^2}{2!}=\frac1{288},\\
 c_3&=b_3+b_1b_2+\frac{b_1^3}{3!}
     =-\frac{139}{51840},\\
 c_4&=b_4+b_1b_3+\frac{b_2^2}{2!}
      +\frac{b_1^2b_2}{2!}+\frac{b_1^4}{4!}
     =-\frac{571}{2488320},\\
 c_5&=b_5+b_2b_3+b_1b_4
      +\frac{b_1^2b_3}{2!}
      +\frac{b_1b_2^2}{2!}
      +\frac{b_1^3b_2}{3!}
      +\frac{b_1^5}{5!}
     =\frac{163879}{209018880}.
\end{align*}
These examples illustrate why the Bell-polynomial notation is not cosmetic:
at order $j$, one must sum over every partition of $j$ using the admissible
odd part sizes.  Formula \cref{eq:c-odd-partitions} performs that bookkeeping
at arbitrary order.

For a fixed pole $k$, multiplication by $\rho/\zeta_k$ gives immediately
\[
 C_{0,k}=\frac{\rho}{\zeta_k},\qquad
 C_{1,k}=\frac{\rho}{12\zeta_k},\qquad
 C_{2,k}=\frac{\rho}{288\zeta_k},\qquad
 C_{3,k}=-\frac{139\rho}{51840\zeta_k},\ldots.
\]
No second recurrence in $k$ is needed.

\section{A proof-oriented transseries dictionary}\label{app:dictionary}
\begin{longtable}{@{}p{0.24\textwidth}p{0.68\textwidth}@{}}
\caption{Objects in the Fubini calculation and their analytic roles.}
\label{tab:dictionary}\\
\toprule
Object & Role\\
\midrule
\endfirsthead
\toprule
Object & Role\\
\midrule
\endhead
$\Fub_n=\sum_jj!\StirlingII{n}{j}$
& Exact ordered-partition count; a Stirling transform of factorials.\\[1mm]
$(2-\e^z)^{-1}$
& Exponential generating function; its meromorphic geometry determines the
  exponential sectors.\\[1mm]
$\zeta_k=\rho+2\pi\ii k$
& Complete pole lattice.  Modulus determines suppression; argument determines
  oscillation.\\[1mm]
$u_k=\rho/\zeta_k$
& Dimensionless sector multiplier.  The normalized exact transmonomial is
  $u_k^{n+1}$.\\[1mm]
$\Lambda_k=\Log(\zeta_k/\rho)$
& Complex action; $u_k^{n+1}=\e^{-n\Lambda_k}\e^{-\Lambda_k}$.\\[1mm]
$T_n=\sum_ku_k^{n+1}$
& Exact, absolutely convergent pole-lattice transseries in factorial
  normalization.\\[1mm]
$\GammaStar(n)$
& Exact common factorial fluctuation in elementary normalization.\\[1mm]
$b_r=B_{r+1}/[r(r+1)]$
& Coefficients of the logarithmic Stirling series.\\[1mm]
$c_j$
& Coefficients after exponentiation; complete Bell polynomials in the $b_r$.\\[1mm]
$C_{j,k}=c_j\rho/\zeta_k$
& Complete mixed coefficient array for algebraic order $j$ and pole sector
  $k$.\\[1mm]
$B_{n,K}$
& Rigorous envelope of all omitted pole pairs beyond $K$.\\[1mm]
$G_J(n)$
& Rigorous scaled-gamma remainder after $J$ algebraic terms.\\[1mm]
finite multi-index $\bm m$
& Records products of nonprincipal pole sectors created by nonlinear
  operations.\\
\bottomrule
\end{longtable}

\section{Continuous-order asymptotics}\label{app:continuous}
For real $\nu>0$, the exact formula \cref{eq:continuous-poles} may be written
\begin{equation}\label{eq:continuous-factorization}
 \mathscr F(\nu)
 =\frac{\sqrt{2\pi\nu}}{2\rho}
  \left(\frac{\nu}{\e\rho}\right)^\nu
  \GammaStar(\nu)
  \sum_{k\in\mathbb Z}u_k^{\nu+1}.
\end{equation}
Here powers use the principal logarithm.  Since $u_{-k}=\overline{u_k}$, the
sum is real for real $\nu$.  Therefore
\begin{align}
 \mathscr F(\nu)
 ={}&\frac{\sqrt{2\pi\nu}}{2\rho}
  \left(\frac{\nu}{\e\rho}\right)^\nu
  \GammaStar(\nu)\nonumber\\
 &\times\left[1+2\sum_{k\ge1}
  q_k^{\nu+1}\cos((\nu+1)\theta_k)\right].
 \label{eq:continuous-real}
\end{align}
The same coefficients $c_j$ and $C_{j,k}$ apply with $n$ replaced by $\nu$.
For complex $\nu$ with $\Re\nu>0$, the bilateral sum remains absolutely
convergent, but conjugate pairing no longer produces a real expression and
the relative dominance of $k$ and $-k$ depends on $\Im\nu$.

\section{Sharper elementary pole-tail estimates}\label{app:tail}
The Hurwitz-zeta bound in \cref{eq:B-hurwitz} is simple and uniform.  Several
refinements are occasionally useful.  Since $q_k$ decreases with $k$,
\begin{equation}\label{eq:tail-first-plus}
 B_{n,K}=2q_{K+1}^{n+1}
 \left[1+\sum_{r\ge1}
  \left(\frac{R_{K+1}}{R_{K+1+r}}\right)^{n+1}\right].
\end{equation}
Thus the first omitted conjugate pair gives the correct exponential scale.
Using $R_k\ge2\pi k$ only in the remaining tail yields
\begin{equation}\label{eq:tail-hybrid}
 B_{n,K}
 \le2q_{K+1}^{n+1}
 +2\left(\frac{\rho}{2\pi}\right)^{n+1}
   \zeta(n+1,K+2).
\end{equation}
This is often considerably sharper than applying the elementary majorant to
the first omitted pair as well.

For $K=0$, a phase-blind two-sided estimate follows from
\cref{eq:global-relative-bound}:
\begin{equation}\label{eq:two-sided-leading}
 1-\frac{\rho^2}{12}
 \le\frac{2\rho^{n+1}\Fub_n}{n!}
 \le1+\frac{\rho^2}{12},
 \qquad n\ge1.
\end{equation}
The lower bound is positive, so logarithms, reciprocals, and arbitrary real
powers of the normalized sequence are globally well-defined for all positive
integer $n$.

\section{Generalized Bernoulli coefficients for multiple poles}
\label{app:norlund}
The generalized Bernoulli numbers in \cref{eq:norlund-def} themselves admit
Bell-polynomial formulae.  Write
\[
 \log\!\left(\frac{t}{\e^t-1}\right)
 =-\frac t2-
 \sum_{m\ge1}\frac{B_{2m}}{2m(2m)!}t^{2m}.
\]
If $\ell_j$ denotes the coefficient of $t^j$ in this logarithm, then
\[
 B_j^{(r)}=
 \BellY_j(1!r\ell_1,2!r\ell_2,\ldots,j!r\ell_j).
\]
Writing
\(
 \log(t/(\e^t-1))=\sum_{m\ge1}\ell_m t^m
\), logarithmic differentiation of \cref{eq:norlund-def} gives the
triangular recurrence
\begin{equation}\label{eq:norlund-recurrence}
 B_0^{(r)}=1,
 \qquad
 B_j^{(r)}=
 r\sum_{m=1}^{j}\binom{j-1}{m-1}
 (m!\ell_m)\,B_{j-m}^{(r)}.
\end{equation}
Here $\ell_1=-1/2$, $\ell_{2m}=-B_{2m}/[2m(2m)!]$, and
$\ell_{2m+1}=0$ for $m\ge1$.  Any preferred standard recurrence for
N\o rlund numbers may be substituted in \cref{eq:higher-fubini-poles}; the
pole formula itself requires only the generating definition
\cref{eq:norlund-def}.

\begin{thebibliography}{99}\sloppy\raggedright

\bibitem{Ahlbach2013}
C.~Ahlbach, J.~Usatine, and N.~Pippenger,
``Barred preferential arrangements,''
\emph{Electronic Journal of Combinatorics} \textbf{20} (2013), no.~2,
Paper P55.

\bibitem{Comtet1974}
L.~Comtet,
\emph{Advanced Combinatorics: The Art of Finite and Infinite Expansions},
D.~Reidel, Dordrecht, 1974.

\bibitem{DLMF}
NIST Digital Library of Mathematical Functions,
``Gamma Function, \S5.11: Asymptotic Expansions,''
\url{https://dlmf.nist.gov/5.11}, accessed September 2026.

\bibitem{FlajoletSedgewick2009}
P.~Flajolet and R.~Sedgewick,
\emph{Analytic Combinatorics}, Cambridge University Press, Cambridge, 2009.

\bibitem{Good1975}
I.~J.~Good,
``The number of orderings of $n$ candidates when ties are permitted,''
\emph{The Fibonacci Quarterly} \textbf{13} (1975), no.~1, 11--18.
\url{https://doi.org/10.1080/00150517.1975.12430678}.

\bibitem{Gross1962}
O.~A.~Gross,
``Preferential arrangements,''
\emph{American Mathematical Monthly} \textbf{69} (1962), no.~1, 4--8.

\bibitem{Hardy1949}
G.~H.~Hardy,
\emph{Divergent Series}, Oxford University Press, Oxford, 1949.

\bibitem{Norlund1924}
N.~E.~N\o rlund,
\emph{Vorlesungen \"uber Differenzenrechnung}, Springer, Berlin, 1924.

\bibitem{OEIS}
OEIS Foundation Inc.,
``The On-Line Encyclopedia of Integer Sequences, A000670: Fubini numbers,''
\url{https://oeis.org/A000670}, accessed September 2026.

\bibitem{Olver1997}
F.~W.~J.~Olver,
\emph{Asymptotics and Special Functions}, AKP Classics, A K Peters,
Wellesley, MA, 1997.

\bibitem{ParisKaminski2001}
R.~B.~Paris and D.~Kaminski,
\emph{Asymptotics and Mellin--Barnes Integrals},
Cambridge University Press, Cambridge, 2001.

\bibitem{Pippenger2010}
N.~Pippenger,
``The hypercube of resistors, asymptotic expansions, and\linebreak[3]
preferential arrangements,''
\emph{Mathematics Magazine} \textbf{83} (2010), no.~5, 331--346.
\href{https://doi.org/10.4169/002557010X529752}{doi:10.4169/002557010X529752}.

\bibitem{ProveItCoefficient}
V.~Reshetnikov,
``Combinatorial Coefficient Calculus,'' research draft in the ProveIt
repository,
\url{https://github.com/VladimirReshetnikov/ProveIt/tree/main/Analysis/FabiusFunction/docs/semi-formalized-research-frontiers/drafts/combinatorial-coefficient-calculus/Combinatorial_Coefficient_Calculus},
accessed September 2026.

\bibitem{ProveItTransseries}
V.~Reshetnikov,
``Series and transseries,'' research-draft collection in the ProveIt
repository,
\url{https://github.com/VladimirReshetnikov/ProveIt/tree/main/Analysis/FabiusFunction/docs/semi-formalized-research-frontiers/drafts/series-and-transseries},
accessed September 2026.

\bibitem{WhittakerWatson1927}
E.~T.~Whittaker and G.~N.~Watson,
\emph{A Course of Modern Analysis}, 4th ed., Cambridge University Press,
Cambridge, 1927.

\end{thebibliography}

\end{document}
